\section{AMC 10 10B 2024}
        \begin{enumerate}
        	\item (\textit{AMC 10 10B 2024}) In a long line of people arranged left to right, the 1013th person from the left is also the 1010th person from the right. How many people are in the line?


 $\textbf{(A) } 2021 \qquad\textbf{(B) } 2022 \qquad\textbf{(C) } 2023 \qquad\textbf{(D) } 2024 \qquad\textbf{(E) } 2025$

 \textbf{Answer: }B



\textbf{Solution for Certain China Testpapers}

If the person is the 1015th from the left, that means there is 1014 people to their left.If the person is the 1010th from the right, that means there is 1009 people to their right.Therefore, there are $1014 + 1 + 1009 = \boxed{\textbf{(D) } 2024}$ people in line.

 



\textbf{Solution 1}

If the person is the 1013th from the left, that means there is 1012 people to their left.If the person is the 1010th from the right, that means there is 1009 people to their right.Therefore, there are $1012 + 1 + 1009 = \boxed{\textbf{(B) } 2022}$ people in line.

 



\textbf{Solution 2}

The person is $1013$rd person from the left is also the $1010$th person from the right, so the same person is counted twice.

 Therefore, there are $1013+1010-1=\boxed{\textbf{(B) } 2022}$ people in line.

 $\boxed{\textbf{(B) } 2022}$




\textbf{Solution 3}

We can look at a smaller case where it is the $4$th person from the left and the $2$nd person from the right. Listing out the people as numbers gives us a list of $1,2,3,4,5.$ We see that the total number of people is the position from the left + the position from the right - 1, which in the case above is $4+2-1=5.$ Plugging in $1013$ and $1010$ gives us $1013+1010-1=\boxed{\textbf{(B) } 2022}$




\textbf{Solution 4}

Since the person is the 1013th person from the left and the 1010th person from the right, there are 1009 people after the 1013th person. 

 Therefore, there are 1013 + 1009 = $\boxed{\textbf{(B) } 2022}$ people in line.

 (Note: Similarly, you can perceive this information as 1012 people to the left of 1010 people and proceed.)

 



\textbf{Solution 5 (Think of a Pattern)}

If you are second going from the left and third going from the right, there are the third person out of four people $p_1,\; p_2,\; p_3,\; p_4$.

 We can simply add $1013+1010$ and subtract $1$ since you are counting yourself twice, giving $1013 + 1009$ = $\boxed{\textbf{(B) } 2022}$



	\item (\textit{AMC 10 10B 2024}) What is $10! - 7! \cdot 6!?$
$\textbf{(A) } -120 \qquad\textbf{(B) } 0 \qquad\textbf{(C) } 120 \qquad\textbf{(D) } 600 \qquad\textbf{(E) } 720$

 \textbf{Answer: }B



\textbf{Solution 1}

$10! = 10 \cdot 9 \cdot 8 \cdot 7! = 720 \cdot 7!$
$6! \cdot 7! = 720 \cdot 7!$

 Therefore, the equation is equal to $720 \cdot 7! - 720 \cdot 7! = \boxed{\textbf{(B) }0}$

 [ONLY FOR CERTAIN CHINESE TESTPAPERS]

 $0 - 5! = \boxed{\textbf{(A) }-120}$




\textbf{Solution 2}

Factoring out $7!$ gives 
 \[7!(10\cdot9\cdot8-1\cdot6!).\] Since $10\cdot9\cdot8=6!=720$, the answer is $\boxed{\text{(B) }0}$ 

 Factoring $6!$ also works, it just makes the expression in the parenthesis a little harder to compute.

 



\textbf{Solution 3}

Note that $10! - 7! \cdot 6!$ must be divisible by $7$, since 
 \[10! - 7! \cdot 6!=7!(10\cdot9\cdot8\cdot6!)=7\cdot(6!\cdot10\cdot9\cdot8\cdot6!)\] and $\boxed{\text{(B) }0}$ is the only option divisible by $7$.

 



\textbf{Solution 4}

$10! - 7! \cdot 6!$ can be split into two parts, $10!$ and $7! \cdot 6!$. We can break the $6!$ into $(2 \cdot 4)(3 \cdot 5 \cdot 6)$The $2 \cdot 4$ part makes $8$, and the $3 \cdot 5 \cdot 6$ part makes $90$, which is $9 \cdot 10$. We still have the 7!, and we can multiply it by $8 \cdot 9 \cdot 10$. This is clearly equivalent to $10!$, so our solution is $10! - 10! =$$\boxed{\text{(B) }0}$.

 



\textbf{Solution 5}

$10! = 3,628,800$, $7! = 5,040$, and $6! = 720$. $5,040 \cdot 720 = 3,628,800$. Therefore, $3,628,800 - 3,628,800 = \boxed{\text{(B) }0}$.

 -pepper2831

 -minor edits by blakezhang

 


	\item (\textit{AMC 10 10B 2024}) For how many integer values of $x$ is $|2x| \leq 7 \pi$?


 $\textbf{(A) } 16 \qquad\textbf{(B) } 17 \qquad\textbf{(C) } 19 \qquad\textbf{(D) } 20 \qquad\textbf{(E) } 21$

 \textbf{Answer: }E



\textbf{Solution for Certain China Testpapers}

Use the fact that $\pi \approx 3.14$, and thus you can get $6\pi \approx 18.84$. We could easily see that the answer is $\{-9,-8,...,8,9\}\implies\boxed{\text{(C) }19}$




\textbf{Solution 1}

$\pi = 3.14159\dots$ is slightly less than $\dfrac{22}{7} = 3.\overline{142857}$. So $7\pi \approx 21.9$The inequality expands to be $-21.9 \le 2x \le 21.9$. We find that $x$ can take the integer values between $-10$ and $10$ inclusive. There are $\boxed{\text{E. }21}$ such values.

 Note that if you did not know whether $\pi$ was greater than or less than $\dfrac{22}{7}$, then you might perform casework. In the case that $\pi > \dfrac{22}{7}$, the valid solutions are between $-11$ and $11$ inclusive: $23$ solutions. Since, $23$ is not an answer choice, we can be confident that $\pi < \dfrac{22}{7}$, and that $\boxed{\text{E. } 21}$ is the correct answer.

 Test advice: If you are in the test and do not know if $\frac{22}{7}$ is bigger or smaller than $\pi$, you can use the extremely sophisticated method of just dividing $\dfrac{22}{7}$ via long division. Once you get to $3.142$ you realise that you don't need to divide further since $\pi = 3.1416$ when rounded to 4 decimal places.Therefore, you do not include $11$ and $-11$ and the answer is 21.

 



\textbf{Solution 2}

$7\pi$ is incredibly close to $22$, but doesn't reach it. This can be both computed by using $\pi\approx3.142\implies7\cdot3.142=21.994<22$ or assumed. Therefore, including both positive and negative values, the answer is $\{-10,-9,...,9,10\}\implies\boxed{\text{(E) }21}$. ~Tacos\_are\_yummy\_1

 


	\item (\textit{AMC 10 10B 2024}) Balls numbered $1, 2, 3, \dots$ are deposited in $5$ bins, labeled $A$, $B$, $C$, $D$, and $E$ using the following procedure. Ball $1$ is deposited in bin $A$, and balls $2$ and $3$ are deposited in bin $B$. The next $3$ balls are deposited in bin $C$, the next $4$ in bin $D$, and so on, cycling back to bin $A$ after balls are deposited in bin $E$. (For example, balls numbered $22, 23, \dots, 28$ are deposited in bin $B$ at step $7$ of this process.) In which bin is ball $2024$ deposited?


 $\textbf{(A) } A \qquad\textbf{(B) } B \qquad\textbf{(C) } C \qquad\textbf{(D) } D \qquad\textbf{(E) } E$

 \textbf{Answer: }D



\textbf{Solution 1}

Consider the triangular array of numbers:
 \[1\]
 \[2, 3\]
 \[4, 5, 6\]
 \[7, 8, 9, 10\]
 \[11, 12, 13, 14, 15\]
 \[\vdots\].

 The numbers in a row congruent to $1 \bmod{5}$ will be in bucket A. Similarly, the numbers in a row congruent to $2, 3, 4, 0 \bmod{5}$ will be in buckets B, C, D, and E respectively. Note that the $n^\text{th}$ row ends with the $n^\text{th}$ triangle number, $\frac{n(n+1)}{2}$.

 We must find values of $n$ that make $\frac{n(n+1)}{2}$ close to $2024$.
 \[\frac{n(n+1)}{2} \approx 2024\]
 \[n(n+1) \approx 4048\]
 \[n^2 \approx 4048\]
 \[n \approx 63\]

 Trying $n = 63$ we find that $\frac{n(n+1)}{2} = 2016$. Since $2016$ will be the last ball in row $63$, ball $2024$ will be in row $64$. Since $64 \equiv 4 \bmod{5}$, ball $2024$ will be placed in bucket $\boxed{\text{D. } D}$.

 



\textbf{Solution 2}


                    \begin{figure}[H]
                        \centering
                        \includegraphics[width=0.5\textwidth]{Solution/AMC_10_10B_2024_4_Solution_Figure_1.jpeg}
                    \end{figure}
                    


	\item (\textit{AMC 10 10B 2024}) In the following expression, Melanie changed some of the plus signs to minus signs:

 \[1+3+5+7+\cdots+97+99\]
When the new expression was evaluated, it was negative. What is the least number of plus signs that Melanie could have changed to minus signs?


 $\textbf{(A) } 14 \qquad\textbf{(B) } 15 \qquad\textbf{(C) } 16 \qquad\textbf{(D) } 17 \qquad\textbf{(E) } 18$

 \textbf{Answer: }B



\textbf{Solution 1}

Recall that the sum of the first $n$ odd numbers is $n^2$. Thus 
 \[1 + 3 + 5 + 7+ \dots + 97 + 99 = 50^2 = 2500.\]

 If we want to minimize the number of sign flips to make the number negative, we must flip the signs corresponding to the values with largest absolute value. This will result in the inequality 
 \[1 + 3 + 5 +\dots + (2n - 3) + (2n - 1) - (2n + 1) - (2n + 3)-\dots - 97 - 99 < 0.\]

 The positive section of the sum will contribute $n^2$, and the negative section will contribute $-(2500-n^2) = (n^2 - 2500)$. The inequality simplifies to
 \[n^2 + (n^2 - 2500) < 0\]
 \[2n^2 < 2500\]
 \[n^2 < 1250\]The greatest positive value of $n$ satisfying the inequality is $n = 35$, corresponding to $35$ positive numbers, and $\boxed{\text{B. } 15}$ negatives.

 ALEX

 



\textbf{Solution 2}

The formula for the sum of all odd positive integers from $1$ to $2n-1$ is 
 \[1+3+5+...+2n-3+2n-1=n^2.\]Therefore, the given sum evaluates to 
 \[99=2(50)-1\implies 50^2=2500.\] Since we're looking for the minimum possible sign changes, we focus on the largest numbers in the set to change to negative.

 If we change the sign of a number $n$ to negative, then the sum decreases by $2n$. Therefore, we're looking for a subset of numbers that add to greater than $2500/2=1250$.

 Now we look at the answer choices.

 $\text{(A) }14$ means that we're changing the signs of the numbers $\{73,75,...,99\}$, and $73+75+...+99=\dfrac{(73+99)\cdot(\frac{99-73}{2}+1)}{2}=1204$. 

 Now this so slightly happens to be less than $1250$, and $\text{(B) } 15$ means we're adding $71$ to the set, too. Since $71>1250-1204=46$, then the answer is $\boxed{\text{(B) }15}$ ~Tacos\_are\_not\_yummy\_1

 



\textbf{Solution 3}

Note that as mentioned above, $1 + 3 + 5... + 2n - 1 = n^{2}$. So, the sum of the first 50 odd numbers is equal to 2500, and we need to find the smallest integer such that $2k^{2} < 50^{2}$. Taking the square root, $k = 25\sqrt{2}$, keeping in mind we need to take the floor $2k^{2}$ needs to be slightly less than 2500. Using an approximation of $\sqrt{2} = 1.4, 25 * \frac{7}{5}$ is 35, so our answer is $\boxed{\text{(B) }15}$~aleyang

 



\textbf{Solution 4 (Simple Algebra)}

As mentioned above, $1 + 3 + 5... + 2n - 1 = n^{2}$. This will be especially important in this solution. The sum of the original expression is 2500, so the combined (positive) sum of the negatives must be greater than half of this, i.e. $1250$. We are looking to minimize the number of negative numbers, so we only convert the last numbers to negatives. Let $99 = o_1 = 100-1$, $97 = o_2 = 100 - 3$, and for any $n$, $o_n = 100-(2n-1)$. Then the sum of all numbers up till $o_n$ is $100n - n^{2}$. This sum needs to be greater than 1250, i.e.:

 $100n-n^{2} > 1250$

 Testing out $n$ from the solution set yields the minimum value of $n=15$. The answer is therefore $\boxed{\text{(B) }15}$.

 


	\item (\textit{AMC 10 10B 2024}) A rectangle has integer side lengths and an area of $2024$. What is the least possible perimeter of the rectangle?


 $\textbf{(A) } 160 \qquad\textbf{(B) } 180 \qquad\textbf{(C) } 222 \qquad\textbf{(D) } 228 \qquad\textbf{(E) } 390$

 \textbf{Answer: }B



\textbf{Solution 1 - Prime Factorization}

We can start by assigning the values x and y for both sides. Here is the equation representing the area:

 
$x \cdot y = 2024$

 Let's write out 2024 fully factorized.

 
$2^3 \cdot 11 \cdot 23$

 Since we know that $x^2 > (x+1)(x-1)$, we want the two closest numbers possible. After some quick analysis, those two numbers are $44$ and $46$. $\\44+46=90$

 Now we multiply by $2$ and get $\boxed{\textbf{(B) }180}.$

 Solution by \href{https://artofproblemsolving.com/wiki/index.php?title=User:IshikaSaini&action=edit&redlink=1}{IshikaSaini}.

 



\textbf{Solution 2 - Squared Numbers Trick}

We know that $x^2 = (x-1)(x+1)+1$ . Recall that $45^2 = 2025$. 

 If I want 1 less than 2025, which is 2024, I can take 1 number higher and 1 number lower from 45, which are 46 and 44. These are the 2 sides of the minimum perimeter because the 2 numbers are closest to each other, which is what we want to get the minimum. 

 Finding the perimeter with $2(46+44)$ we get $\boxed{\textbf{(B) }180}.$

 Solution by ~Taha Jazaeri

 Note: The square of any positive integer with units digit $5$, written in the format $10x + 5$ where $x$ is a positive integer, is equal to $100(x)(x+1)+25$. 

 



\textbf{Solution 3 - AM-GM Inequality}

Denote the numbers as $x, \frac{2024}{x}$. We know that per AM-GM, $x+\frac{2024}{x} \geq 2\sqrt{2024}$, but since $2\sqrt{2025} = 90$, $2\sqrt{2024}$ must be slightly less than 90, so $2x + 2\frac{2024}{x}$ must be slightly less than 180, eliminating A as a possible answer choice. Proceed with the following solutions above to get 44 and 46, which is $\boxed{\textbf{(B) }180}.$

 -aleyang

 



\textbf{Solution 4 - Difference of Squares}

Note that the year 2025 is a perfect square. The beauty of this year shines down to 2024, which is 1 year lower. $45^2-1^2=2024=(45-1)(45+1)$

 Knowing this is the closest possible we can get to a square with only integer factors of 2024, these two are our happy numbers! Add them up to get:

 



\textbf{Solution 5 - Get Lucky}

Note: This is what I did.

 $\sqrt{2025}=45$

 Assuming it's a square,

 $45\cdot4=\boxed{\textbf{180}}$




\textbf{Solution 6 (Better Prime Factorization)}

The smallest perimeter is achieved if the two numbers are as close as possible to one another. The prime factorization of 2024 is \( 2^3 \cdot 11 \cdot 23 \). Notice that 11 is smaller than 23 and there are three 2's to be distributed between 11 and 23. We give 11 \( 2^2 \) and 23 \( 2^1 \). This guarantee's the smallest perimeter. The side length of the rectangle is \( 44 \times 46 \). The perimeter of a rectangle is \( 2l + 2w \), where \( l = 46 \) and \( w = 44 \), so therefore the perimeter is \( 2(46) + 2(44) = 180 \). Therefore the answer is $\boxed{\textbf{(B)} 180}$.

 


	\item (\textit{AMC 10 10B 2024}) What is the remainder when $7^{2024}+7^{2025}+7^{2026}$ is divided by $19$?


 $\textbf{(A) } 0 \qquad\textbf{(B) } 1 \qquad\textbf{(C) } 7 \qquad\textbf{(D) } 11 \qquad\textbf{(E) } 18$

 \textbf{Answer: }A



\textbf{Solution 1}

We can factor the expression as

 
 \[7^{2024} (1 + 7 + 7^2) = 7^{2024} (57).\]

 Note that $57=19\cdot3$, this expression is actually divisible by 19. The answer is $\boxed{\textbf{(A) } 0}$.

 



\textbf{Solution 2}

If you failed to realize that the expression can be factored, you might also apply modular arithmetic to solve the problem.

 Since $7^3\equiv1\pmod{19}$, the powers of $7$ repeat every three terms:

 
 \[7^1\equiv7\pmod{19}\]
 \[7^2\equiv11\pmod{19}\]
 \[7^3\equiv1\pmod{19}\]

 The fact that $2024\equiv2\pmod3$, $2025\equiv0\pmod3$, and $2026\equiv1\pmod3$ implies that $7^{2024}+7^{2025}+7^{2026}\equiv11+1+7\equiv19 \equiv0\pmod{19}$.

 



\textbf{Solution 3}

We start the same as solution 2, and find that:

 
 \[7^1\equiv7\pmod{19}\]
 \[7^2\equiv11\pmod{19}\]
 \[7^3\equiv1\pmod{19}\]

 We know that for $2024$, $2025$, and $2026$, because there are three terms, we can just add them up. $1 + 7 + 11 = 19$, which is $0$ mod $19$. 

 


 



\textbf{Solution 4 (Given more advanced knowledge)}

By Fermat's Little Theorem (FLT), we know that 
 \[7^{18}\equiv1\pmod{19}\] Then its order must divide $18$. Trying simple values we try and succeed: 
 \[7^3\equiv1\pmod{19}\]

 
So the expression is equivalent to $1+7+49\pmod{19}$, which gives $\boxed{\textbf{(A) } 0}$ when divided by 19.

 



\textbf{Solution 5 (Basic Mod Almost same as sol one)}

We have \begin{align*}7^{2024} \cdot (1+7+49) &\equiv 0 \pmod{19}\\ 57 &\equiv 0 \pmod{19},\end{align*}
 so the answer is $\boxed{\textbf{(A) } 0}$. 

 


	\item (\textit{AMC 10 10B 2024}) Let $N$ be the product of all the positive integer divisors of $42$. What is the units digit
of $N$?


 $\textbf{(A) } 0\qquad\textbf{(B) } 2\qquad\textbf{(C) } 4\qquad\textbf{(D) } 6\qquad\textbf{(E) } 8$

 \textbf{Answer: }D



\textbf{Solution 1}

The factors of $42$ are $1, 2, 3, 6, 7, 14, 21, 42$. Multiply the unit digits to get $\boxed{\textbf{(D) } 6}$




\textbf{Solution 2}

The product of the factors of a number $n$ is $n^\frac{\tau(n)}{2}$, where $\tau(n)$ is the number of positive divisors of $n$. We see that $42 = 2^1 \cdot 3^1 \cdot 7^1$ has $(1+1)(1+1)(1+1) = 8$ factors, so the product of the divisors of $42$ is

 
 \[42^{\frac{8}{2}} \equiv 42^4 \equiv 2^4 \equiv 16 \equiv \boxed{\textbf{(D) } 6} \pmod{10}.\]




\textbf{Solution 3}

The factors are $1, 2, 3, 6, 7, 14, 21, 42$. We can take out the 1 giving us $2, 3, 6, 7, 14, 21, 42$. We can then take the unit digit of each number because the unit digit of the product is dependent on the product of the unit digits in the factors. So the new list is $2, 3, 6, 7, 4, 1, 2$. Taking out the 1 again we get $2, 3, 6, 7, 4, 2$. Multiplying the first 3 numbers gives us $36$ which we take the unit digit off. Taking the final 3 digits we get $56$ which we take the unit digit off again. This leaves us with 2 numbers which are $6$ and $6$, the product of them is 36 and the unit digit of this is $\boxed{\textbf{(D) } 6}$



	\item (\textit{AMC 10 10B 2024}) Real numbers $a, b,$ and $c$ have arithmetic mean $0$. The arithmetic mean of $a^2, b^2,$ and $c^2$ is $10$. What is the arithmetic mean of $ab, ac,$ and $bc$?


 $\textbf{(A) } -5 \qquad\textbf{(B) } -\dfrac{10}{3} \qquad\textbf{(C) } -\dfrac{10}{9} \qquad\textbf{(D) } 0 \qquad\textbf{(E) } \dfrac{10}{9}$

 \textbf{Answer: }A



\textbf{Solution 1}

If $\frac{a+b+c}{3} = 0$, that means $a+b+c=0$, and $(a+b+c)^2=0$. Expanding that gives 
 \[(a+b+c)^2=a^2+b^2+c^2+2ab+2ac+2bc=0\] If $\frac{a^2+b^2+c^2}{3} = 10$, then $a^2+b^2+c^2=30$. Thus, we have 
 \[30 + 2ab + 2ac + 2bc = 0\] Basic arithmetic will give you that $ab + bc + ac = -15$. To find the arithmetic mean, divide that by 3, so $\frac{ab + bc + ac}{3} = \boxed{\textbf{(A) }-5}$




\textbf{Solution 3}

Assume that $a = 0$ and $b = -c$. Since we know that the arithmetic mean of the three numbers is $10$, this means $a^2+b^2+c^2=30$. Using this equation, $b^{2} + b^{2} = 30$, so $b^2 = 15$. Observe that taking the positive or negative root won't matter as $c$ will be the opposite. If we let $b = \sqrt{15}$ and $c = -\sqrt{15}$, $ab = 0\times\sqrt{15} = 0$, $ac = 0\times-\sqrt{15} = 0$, and $bc = \sqrt{15}\times-\sqrt{15} = -15$, so $ab+ac+bc=-15$. Doing $\frac{-15}{3}$ to get the arithmetic mean will give us $\boxed{\textbf{(A)}\ -5}$.

 -aleyang

 



\textbf{Solution 4 (interesting property)}

From the definition of arithmetic means and some simple manipulation, we have $a+b+c=0$ and $a^2+b^2+c^2=30$. Let $ab+ac+bc=k$, we have: 

 
 \[k=ab+ac+bc=(a+b)c+ab\]

 Substituting $c=-(a+b)$ into the expression, we have

 
 \[-c^2+ab=k\]

 Adding twice the equation above to $a^2+b^2+c^2=30$, we have:

 
 \[a^2+b^2+c^2-2c^2+2ab=30+2k\]

 \[a^2+b^2+2ab-c^2=30+2k\]

 \[(a+b)^2-c^2=30+2k\]

 \[(a+b+c)(a+b-c)=30+2k\]

 But we know that $a+b+c=0$. So $30+2k=0$ and $k=-15$. Therefore the arithmetic mean of $ab,ac,bc$ is $\frac{k}{3}=-5$. So the answer is $\boxed{\textbf{(A)}\ -5}$



	\item (\textit{AMC 10 10B 2024}) Quadrilateral $ABCD$ is a parallelogram, and $E$ is the midpoint of the side $\overline{AD}$. Let $F$ be the intersection of lines $EB$ and $AC$. What is the ratio of the area of quadrilateral $CDEF$ to the area of $\triangle CFB$?


 $\textbf{(A) } 5:4 \qquad\textbf{(B) } 4:3 \qquad\textbf{(C) } 3:2 \qquad\textbf{(D) } 5:3 \qquad\textbf{(E) } 2:1$

 \textbf{Answer: }A



\textbf{Solution 1}

Let $AB = CD$ have length $b$ and let the altitude of the parallelogram perpendicular to $\overline{AD}$ have length $h$.

 The area of the parallelogram is $bh$ and the area of $\triangle ABE$ equals $\frac{(b/2)(h)}{2} = \frac{bh}{4}$. Thus, the area of quadrilateral $BCDE$ is $bh - \frac{bh}{4} = \frac{3bh}{4}$.

 We have from $AA$ that $\triangle CBF \sim \triangle AEF$. Also, $CB/AE = 2$, so the length of the altitude of $\triangle CBF$ from $F$ is twice that of $\triangle AEF$. This means that the altitude of $\triangle CBF$ is $2h/3$, so the area of $\triangle CBF$ is $\frac{(b)(2h/3)}{2} = \frac{bh}{3}$.

 Then, the area of quadrilateral $CDEF$ equals the area of $BCDE$ minus that of $\triangle CBF$, which is $\frac{3bh}{4} - \frac{bh}{3} = \frac{5bh}{12}$. Finally, the ratio of the area of $CDEF$ to the area of triangle $CFB$ is $\frac{\frac{5bh}{12}}{\frac{bh}{3}} = \frac{\frac{5}{12}}{\frac{1}{3}} = \frac{5}{4}$, so the answer is $\boxed{\textbf{(A) } 5:4}$.

 
                    \begin{figure}[H]
                        \centering
                        \includegraphics[width=0.5\textwidth]{Solution/AMC_10_10B_2024_10_Solution_Figure_1.png}
                    \end{figure}
                    



\textbf{Solution 2}

Let $[AFE]=1$. Since $\triangle AFE\sim\triangle CFB$ with a scale factor of $2$, $[CFB]=4$. The scale factor of $2$ also means that $\dfrac{AF}{FC}=\dfrac{1}{2}$, therefore since $\triangle BCF$ and $\triangle BFA$ have the same height, $[BFA]=2$. Since $ABCD$ is a parallelogram, 
 \[[BCA]=[DAC]\implies4+2=1+[CDEF]\implies [CDEF]=5\implies\boxed{\text{(A) }5:4}\]

 \textbf{vladimir.shelomovskii@gmail.com, vvsss}

 



\textbf{Solution 3 (Techniques)}

We assert that $ABCD$ is a square of side length $6$. Notice that $\triangle AFE\sim\triangle CFB$ with a scale factor of $2$. Since the area of $\triangle ABC$ is $18 \implies$ the area of $\triangle CFB$ is $12$, so the area of $\triangle AFE$ is $3$. Thus the area of $CDEF$ is $18-3=15$, and we conclude that the answer is $\frac{15}{12}\implies\boxed{\text{(A) }5:4}$




\textbf{Solution 4}

Let $ABCE$ be a square with side length $1$, to assist with calculations. We can put this on the coordinate plane with the points $D = (0,0)$, $C = (1, 0)$, $B = (1, 1)$, and $A = (0, 1)$. We have $E = (0, 0.5)$. Therefore, the line $EB$ has slope $0.5$ and y-intercept $0.5$. The equation of the line is then $y = 0.5x + 0.5$. The equation of line $AC$ is $y = -x + 1$. The intersection is when the lines are equal to each other, so we solve the equation. $0.5x + 0.5 = -x + 1$, so $x = \frac{1}{3}$. Therefore, plugging it into the equation, we get $y= \frac{2}{3}$. Using the shoelace theorem, we get the area of $CDEF$ to be $\frac{5}{12}$ and the area of $CFB$ to be $\frac{1}{3}$, so our ratio is $\frac{\frac{5}{12}}{\frac{1}{3}} = \boxed{(A) 5:4}$ ~idk12345678

 



\textbf{Solution 5 (wlog)}

Let $ABCE$ be a square with side length $2$. We see that $\triangle AFE \sim \triangle CFB$ by a Scale factor of $2$. Let the altitude of $\triangle AFE$ and altitude of $\triangle CFB$ be $h$ and $2h$, respectively. We know that $h+2h$ is equal to $2$, as the height of the square is $2$. Solving this equation, we get that $h = \frac{2}3.$ This means $[\triangle CFB] = \frac{4}3,$ we can also calculate the area of $\triangle ABE$. Adding the area we of $\triangle CFB$ and $\triangle ABE$ we get $\frac{7}3.$ We can then subtract this from the total area of the square: $4$, this gives us $\frac{5}3$ for the area of quadrilateral $CFED.$ Then we can compute the ratio which is equal to $\boxed{\textbf{(A) } 5:4}.$

 (why does the $\LaTeX$ always look so bugged.)

 



\textbf{Solution 6 (barycentrics)}

$A=(1,0,0), B=(0,1,0), C=(0,0,1), D=(1,-1,1)$. Since $E$ is the midpoint of $\overline{AD}$, $E=(1,-0.5,0.5)$. The equation of $\overline{EB}$ is:
 \[0 = \begin{vmatrix} x & y & z \\ 1 & -0.5 & 0.5 \\ 0 & 1 & 0 \end{vmatrix}\]The equation of $\overline{AC}$ is:
 \[0 = \begin{vmatrix} x & y & z \\ 1 & 0 & 0 \\ 0 & 0 & 1 \end{vmatrix}\]We also know that $x+y+z=1$. To find the intersection, we can solve the system of equations. Solving, we get $x=2/3,y=0,z=1/3$. Therefore, $F=\left(\frac{2}{3}, 0, \frac{1}{3}\right)$. Using barycentric area formula,
 \[\frac{[CFB]}{[ABC]} =  \begin{vmatrix} 0 & 0 & 1 \\ 2/3 & 0 & 1/3 \\ 0 & 1 & 0 \end{vmatrix} =\frac{2}{3}\]
 \[\frac{[CDEF]}{[ABC]} =  \begin{vmatrix} 0 & 0 & 1 \\ 1 & -0.5 & 0.5 \\ 2/3 & 0 & 1/3 \end{vmatrix} + \begin{vmatrix} 0 & 0 & 1 \\ 1 & -1 & 1 \\ 1 & -0.5 & 0.5 \end{vmatrix} =\frac{5}{6}\]$\frac{[CDEF]}{[CFB]}=\frac{\frac{5}{6}}{\frac{2}{3}}=\boxed{\textbf{(A) } 5:4}$



	\item (\textit{AMC 10 10B 2024}) In the figure below $WXYZ$ is a rectangle with $WX=4$ and $WZ=8$. Point $M$ lies on $\overline{XY}$, point $A$ lies on $\overline{YZ}$, and $\angle WMA$ is a right angle. The areas of $\triangle WXM$ and $\triangle W A Z$ are equal. What is the area of $\triangle WMA$?


 
                    \begin{figure}[H]
                        \centering
                        \includegraphics[width=0.5\textwidth]{Problem/AMC_10_10B_2024_11_Problem_Figure_1.png}
                    \end{figure}
                    
$\textbf{(A) }13 \qquad \textbf{(B) }14 \qquad \textbf{(C) }15 \qquad \textbf{(D) }16 \qquad \textbf{(E) }17 \qquad$

 \textbf{Answer: }C



\textbf{Solution 1}

We know that $WX = 4$, $WZ = 8$, so $YZ = 4$ and $YX = 8$. Since $\angle WMA = 90^\circ$, triangles $WXM$ and $MYA$ are similar. Therefore, $\frac{WX}{MY} = \frac{XM}{YA}$, which gives $\frac{4}{8 - XM} = \frac{XM}{4 - ZA}$. We also know that the areas of triangles $WXM$ and $WAZ$ are equal, so $WX \cdot XM = WZ \cdot ZA$, which implies $4 \cdot XM = 8 \cdot ZA$. Substituting this into the previous equation, we get $\frac{4}{8 - 2ZA} = \frac{2ZA}{4 - ZA}$, yielding $ZA = 1$ and $XM = 2$. Thus,

 
 \[\triangle WMA = 4 \cdot 8 - \frac{4 \cdot 2}{2} - \frac{8 \cdot 1}{2} - \frac{6 \cdot 3}{2} = \boxed{\textbf{(C) }15}\]




\textbf{Solution 2}

Let $XM=b$, $ZA = a$, $4\cdot b= 8\cdot a$, $b = 2a$, 
 \[WM^2 + AM^2 = AW^2\]
 \[(b^2+4^2) + (4-a)^2 + (8-b)^2  = (a^2 + 8^2)\]$a=1$, $b=2$ , 
 \[\triangle WMA = area(WXYZ) - \triangle WZA- \triangle WXM- \triangle MYA = 32 - 4-4-9 = \boxed{\textbf{(C) }15}\]




\textbf{Solution 3 (Pythagorean Theorem)}

Assign ZA as $x$, then AY as $4 - x$. Assign XM as $y$ and MY as $8 - y$. Since triangles WXM and WZA are equal in area, we can say $8x = 4y$, so $y = 2x$. Then, therefore, XM is $2x$ and MY has length $8 - 2x$. We can use the Pythagorean theorem to find WM, which is actually $\sqrt{(2x)^2 + 4^2)} = \sqrt{4x^2 + 16}$. We don't factor it yet - we are going to find $x$ again using the Pythagorean Theorem. Similarly, finding MA is just the square root of the squares of AY and MY individually, or $\sqrt{(8 - 2x)^2 + (4 - x)^2} = \sqrt{64 - 32x + 4x^2 + 16 - 8x + x^2} = \sqrt{5x^2 - 40x + 80}$. Then simply, WA is really $\sqrt{x^2 + 64}$.

 Now we have the three sides of the right triangle: $\sqrt{4x^2 + 16}$, $\sqrt{5x^2 - 40x + 80}$, and $\sqrt{x^2 + 64}$. Per the Pythagorean theorem again, we can see $(4x^2 + 16) + (5x^2 - 40x + 80) = (x^2 + 64)$. Combining like terms gives us $8x^2 - 40x + 32 = 0$, then dividing by 8 gives $x^2 - 5x + 4 = 0$. As this elementary and well-known quadratic gives us the roots of $1$ and $4$, we can see it is a bit weird to have $x = 4$, as then point Z is point A. So we'll assume $x = 1$. We have two legs of the triangle by plugging in the sides with x in them, given that $x = 1$: $\sqrt{20}$ and $\sqrt{45}$. We should know that $20 \cdot 45 = 900$, and $\sqrt{900} = 30.$ Dividing by 2 reveals us our answer: $\boxed{\textbf{(C) }15}$

 note: to further prove that 1 is the correct root, note that setting A at Z makes WZA have area 0. Therefore, WXM also has to have an area of 0. This makes it impossible for angle WMA to be a right angle, as WMA would simply become a right triangle with half the area of WXYZ, making MWA the right angle.

 



\textbf{Solution 4 (Similar Triangles)}

We are given $WX = 4$, $WZ = 8$. $\triangle$ WXM and $\triangle$ MYA have equal area, so let $XM = 2x$ and $ZA = x$. $MY = 8-2x$ and $AY = 4-x$. From this, we can conclude that $\frac{MY}{AY} = \frac{8-2x}{4-x} = \frac{2}{1}$

 Since $WM$ intersects parallel lines $WZ$ and $XY$, $\angle ZWM = \angle WMZ$. $\angle ZWM + \angle MWX = 90^\circ$, so $180^\circ - 90^\circ = \angle WMZ + \angle AMY$. Thus, $\angle MWX = \angle AMY$ and $\Delta WXM$ ~ $\Delta MYA$ due to AA Similarity. 

 Corresponding sides of similar triangles are proportional, so $\frac{WX}{XM} = \frac{MY}{AY}$ or $\frac{4}{2x} = \frac{2}{1}$. It is clear that $2x = 2$, and $x = 1$. Now, all we have to do is subtract the area of the rectangle by each of the three triangles.

 
$\Delta WMA$ = $8 \cdot 4$ - $\left(\frac{1}{2} \cdot 4 \cdot 2 \right)$ - $\left(\frac{1}{2} \cdot 8 \cdot 1 \right)$ - $\left(\frac{1}{2} \cdot 6 \cdot 3 \right)$
$\Delta WMA$ = $32 - 4 - 4 - 9$
$\Delta WMA$ = $\boxed{\textbf{(C) }15}$




\textbf{Solution 5 (Desperate)}

Taking a look at the diagram, points $M$ and $A$ seem to divide $\overline{XY}$ and $\overline{YZ}$ into $1:3$ ratios. With this assumption in mind, $XM=2$, $MY=6$, $YA=3$, and $AZ=1$. This means that the areas of $\triangle WXM$ and $\triangle WZA$ are $\frac{4(2)}{2}=4$ and $\frac{8(1)}{2}=4$ respectively, and the area of $\triangle MAY$ is $\frac{6(3)}{2}=9$. The area of rectangle $WXYZ=4(8)=32$, so subtracting the areas of triangles $WXM$, $WZA$, and $MAY$, we get that the area of $\triangle WMA=32-4-4-9=\boxed{\textbf{(C) 15}}$

 \textbf{China Test Solution (Finding $\triangle MAY$)} From solution 3, instead of finding the area of $\triangle WMA$, we instead find  the area of $\triangle MAY$. Then $x = 1$ then we have $MA = 8 - 2x = 6$. Again, since $AY = 4 - x$, then $AY = 4 - 1 = 3.$ The area of a triangle with legs $3$ and $6$ is $\frac{3 * 6}{2} = \boxed{9}$.

 -minor edits by fireball9746

 can someone pls explain why $\triangle MAY$ has the same area as $\triangle WMA$

 \textasciicircum{} To the person above:Under the problem, there’s a note that says some tests in China asked for the area of $\triangle MAY$. - trevian1

 



\textbf{Solution 6 (cheese)}

Of course I don't endorse this method, but...
                    \begin{figure}[H]
                        \centering
                        \includegraphics[width=0.5\textwidth]{Solution/AMC_10_10B_2024_11_Solution_Figure_1.png}
                    \end{figure}
                    

 By looking at partitioned rectangles in the diagram, it's clear that the area of triangle $WMA$ is very close to half the area of the overall rectangle $WXYZ$. So $\boxed{\textbf{(C) 15}}$ is our best bet.

 Obvious note:

 The answer is less than half of the overall area of rectangle $WXYZ$, as triangle $AWZ$ clearly sticks out past rectangle $ACBZ$ into rectangle $XMBW$ by the smallest margin, so therefore, the next smallest answer past (C) 16, is $\boxed{\textbf{(C) 15}}$.

 


	\item (\textit{AMC 10 10B 2024}) A group of $100$ students from different countries meet at a mathematics competition.
Each student speaks the same number of languages, and, for every pair of
students $A$ and $B$, student $A$ speaks some language that student $B$ does not speak,
and student $B$ speaks some language that student $A$ does not speak. What is the
least possible total number of languages spoken by all the students?


 $\textbf{(A) } 9 \qquad\textbf{(B) } 10 \qquad\textbf{(C) } 12 \qquad\textbf{(D) } 51 \qquad\textbf{(E) } 100$

 \textbf{Answer: }A



\textbf{Solution 1 (Pure Logic)}

We think of this problem like boxes. First start with 9. We see that we can arrange the groups of people into the 9 boxes. We take 9 people of different languages and arrange them in each of the boxes. This means we have $100 - 9 \times 9 = 19$ people remaining. We then take 18 people of different language (as they can be put in a pair and still qualify) and put them in the boxes. This gives us 1 person left over. Now, this is where MAA wants you to choose 10, but one can quickly see that the last person can be put into any group of three (except a group that includes their language) and still qualify whilst maintaining all the constraints. Therefore our answer is $\boxed{\textbf{(A) } 9}$.

 This solution is a basic introduction/summary of the \href{https://artofproblemsolving.com/wiki/index.php?title=Pigeonhole\_Principle}{Pigeonhole Principle}

 



\textbf{Solution 2 (Clever intuition)}

There are 100 students at the mathematics convention.

 For any pair of students $AB$ we see that $AB \in \{X,Y\}$ with WLOG, $(X, Y)$ being a unique language individual to either $A$ or $B$.

 In other words, either student $A$ or student $B$ received language $X$ or language $Y$, in which the order they receive their language is unimportant. 

 Given this the number of languages must divide the number of students evenly in order to achieve unique $X,Y$ pairs that satisfy the conditions stated in the problem.

 The question now is "What is the number of languages that divide the number of students?", which is equivalent to asking "How many numbers divide 100?", in which we use the number of factors theorem of the prime factorization of 100, being $2^2 \times 5^2$, and the number of factors being $(2+1)(2+1) = (3)(3) = 9$.

 Therefore our answer is $\boxed{\textbf{(A) } 9}$.

 



\textbf{Solution 3}

Let's say we have some number of languages. Then each student will speak some amount of those languages, and no two people can have the same combination of languages or else the conditions will no longer be satisfied. Notice that $\dbinom94=126\geq100$. So each of the $100$ students can speak some $4$ of the $9$ languages. Thus, $\boxed{9}$ is our answer.

 



\textbf{Solution 4 (Rigorous)}

Let $n$ be the total number of languages spoken by all the students, and let $k$ be the number of languages that each student speaks.Since all students speak the same number of languages, the condition given in the question can be modified as -

 "If $N$ is the $n$-element subset containing the union of the languages spoken by all students, each of the 100 students speaks a different (unique) $k$-element subset of $N$ combination of languages."

 In more mathematical terms, this means ${n \choose k} \ge 100$. (Using \href{https://artofproblemsolving.com/wiki/index.php?title=Pigeonhole\_Principle}{PHP})

 Because we need the least value of $n$, $k$ must be the closest integer to $\frac{n}{2}$
$\Rightarrow {n \choose \frac{n}{2}} \ge 100$ for $n = 2m$, ${n \choose \frac{n-1}{2}} \ge 100$ for $n=2m-1$, $m \in \mathbb{N}$

 In the context of this question, it will be fastest to just substitute $\normalsize n$ to each of the option values to find the least value possible as $\boxed{\textbf{(A) } 9}$

 For a more rigorous proof, where options are not an option, one can manually count up from $n=1$ as $100$ is a small number. 

 


	\item (\textit{AMC 10 10B 2024}) Positive integers $x$ and $y$ satisfy the equation $\sqrt{x} + \sqrt{y} = \sqrt{1183}$. What is the minimum possible value of $x+y$?


 $\textbf{(A) } 585 \qquad\textbf{(B) } 595 \qquad\textbf{(C) } 623 \qquad\textbf{(D) } 700 \qquad\textbf{(E) } 791$

 \textbf{Answer: }B



\textbf{Solution 1}

Note that $\sqrt{1183}=13\sqrt7$. Since $x$ and $y$ are positive integers, and $\sqrt{x}+\sqrt{y}=\sqrt{1183}$ we can represent each value of $\sqrt{x}$ and $\sqrt{y}$ as the product of a positive integer and $\sqrt7$. Let's say that $\sqrt{x}=m\sqrt7$ and $\sqrt{y}=n\sqrt7$, where $m$ and $n$ are positive integers. This implies that 
 \[x+y=(\sqrt{x})^2+(\sqrt{y})^2=7m^2+7n^2=7(m^2+n^2)\] and that $m+n=13$. WLOG, assume that ${m}\geq{n}$. It is not hard to see that $x+y$ reaches its minimum when $m^2+n^2$ reaches its minimum. We now apply algebraic manipulation to get that 
 \[m^2+n^2=(m+n)^2-2mn\] Since $m+n$ is determined, we now want $mn$ to reach its maximum. Since $m$ and $n$ are positive integers, we can use the AM-GM inequality to get that: $\frac{m+n}{2}\geq{\sqrt{mn}}$. When $mn$ reaches its maximum, $\frac{m+n}{2}={\sqrt{mn}}$. This implies that $m=n=\frac{13}{2}$. However, this is not possible since $m$ and $n$ and integers. Under this constraint, we can see that $mn$ reaches its maximum when $m=7$ and $n=6$. Therefore, the minimum possible value of $x+y$ is $7(m^2+n^2)=7(7^2+6^2)=\boxed{\textbf{(B)}595}$

 A similar method is to take $y=1183-26\sqrt{7x}-x^2$, then noting $x=7a^2$ and bashing to find the value of a where x is closest to y.

 



\textbf{Solution 2 (Guessing & Answer Choices)}

Set $x=y$, giving the minimum possible values. The given equation becomes 
 \[\sqrt{x}=\sqrt{y}=\dfrac{\sqrt{1183}}{2}\implies x=y=\dfrac{1183}{4}.\]This means that 
 \[x+y=\dfrac{1183}{2}=591.5.\]Since this is closest to answer choice $\text{(B)}$, the answer is $\boxed{\textbf{(B) }595}$ ~Neoronean ~Tacos\_are\_yummy\_1 (latex)

 Note:If using a solution similar to this one it is recommended to still find valid x and y that add to 595.

 


 



\textbf{Solution 3 (quick, no guessing but still answer choices)}

Square both sides of the equation to get $2\cdot \sqrt{xy} + (x + y) = 1183.$ We can plug in possible values for $x + y$ based on the answer choices. If we do $585,$ we get $\sqrt{xy} = 299 \implies xy = 299^2.$ So $y = \frac{299^2}{x}.$ Then we have$x^2 + 299^2 = 585x,$ which has no real solutions for $x$ since the discriminant is negative. Following the same process for $595$ we obtain the equation $x^2 + 294^2 = 595x,$ which can be factored as $(x - 252)(x - 343) = 0.$ So the answer must be $\boxed{\textbf{(B) }595}$.

 



\textbf{Solution 4 (Calculus)}

Since $\sqrt{x}+\sqrt{y}=\sqrt{1183}$, we can express $y$ in terms of $x$: $y = x+1183-2\sqrt{1183x}$. $x+y = 2x+1183-2\sqrt{1183x}$. Taking the derivative of this, we get $2-2 \cdot \frac{1}{2\sqrt{1183x}} \cdot 1183$. We find the critical points by setting this expression to equate $0$. Simplifying, we get $2 = \frac{1183}{\sqrt{1183x}}$, and then we get $4 \cdot 1183x = 1183^2$. Dividing by $1183$, we get $4x=1183$, so a critical point is $x=\frac{1183}{4}$. Now, we take the second derivative, we get $\frac{13\sqrt{7x}}{2x^2}$. When $x=\frac{1183}{4}$, the second derivative is positive, meaning the minimum value is obtained at $x=\frac{1183}{4}$. This value is $295.75$. However, the problem asks for positive integers $x$ and $y$. We want $x$ to be as close to $295.75$ as possible. We also want $y$ to be an integer, so $x+1183-2\sqrt{1183x}$ should be an integer. Notice how since $x$ is an integer, we just want $\sqrt{1183x}$ to be an integer. Prime factorizing $1183$, we get $1183=7 \cdot 13^2$. Therefore, $x$ must contain an odd amount of $7$'s and an even amount of any other prime factor. We quickly find that $2^2 \cdot 3^2 \cdot 7 = 252$ is the optimal solution for $x$, and plugging in, we get $y=252 + 1183 - 2 \cdot 2 \cdot 3 \cdot 7 \cdot 13 = 343$. We find $x+y=595$, so the answer is choice $\boxed{\text{(B) 595}}$.

 Note:It might not be immediately clear why $x=252$ is the optimal solution. The following explains why $x=252$ is optimal. 

 $x$ must have an odd amount of $7$'s and an even amount of any other prime factor. This means that $x$ can be expressed as $7^n \cdot a$ where $a$ is a perfect square. We can simplify this to $7 \cdot 7^{n-1} \cdot a = 7b$, where $b$ is a perfect square, since $n-1$ is even. This means that $7b$ is as close to $295.75$ as possible. This means that $b \approx \frac{295.75}{7}$. $\frac{295.75}{7}$ is approximately $42$, so $b \approx 42$. Since $b$ is a perfect square, we have two candidates for $b$: $36$ and $49$. If $b=36$, $x=252$, and if $b=49$, $x=343$. To find the optimal value for $x$, pick the value closest to $295.75$, and we find that $\left |295.75-252 \right | < \left |295.75 - 343 \right|$, so $x=252$ is optimal. From now, proceed in the same way as the original solution.

 



\textbf{Solution 5 (cheese)}

$\sqrt{1183}=\sqrt{169\cdot7}=13\sqrt7$.

 Since $\sqrt{x}+\sqrt{y}=13\sqrt7$, let $\sqrt{x}=6\sqrt7$ and $\sqrt{y}=7\sqrt7$.

 Then $\sqrt{x}=\sqrt{252}$ because you take the outside $7$ and put it in. Similarly, $\sqrt{y}=\sqrt{343}$.

 Therefore, the insides sum to $x+y=252+343=595$.

 $\boxed{\text{(B) 595}}$.

 - AskNamanYT

 


	\item (\textit{AMC 10 10B 2024}) A dartboard is the region $B$ in the coordinate plane consisting of points $(x,y)$ such
that $|x| + |y| \le 8$ . A target $T$ is the region where $(x^2 + y^2 - 25)^2 \le 49.$ A dart is
thrown and lands at a random point in $B$. The probability that the dart lands in $T$ can be
expressed as $\frac{m}{n} \cdot \pi,$ where $m$ and $n$ are relatively prime positive integers. What
is $m + n?$
$\textbf{(A) } 39 \qquad\textbf{(B) } 71 \qquad\textbf{(C) } 73 \qquad\textbf{(D) } 75 \qquad\textbf{(E) } 135$

 \textbf{Answer: }B



\textbf{Solution 1}

Inequalities of the form $|x|+|y| \le 8$ are well-known and correspond to a square in space with centre at origin and vertices at $(8, 0)$, $(-8, 0)$, $(0, 8)$, $(0, -8)$.The diagonal length of this square is clearly $16$, so it has an area of 
 \[\frac{1}{2} \cdot 16 \cdot 16 = 128\]Now,
 \[(x^2 + y^2 - 25)^2 \le 49\]Converting to polar form,
 \[r^2 - 25 \le 7 \implies r \le \sqrt{32},\]and
 \[r^2 - 25 \ge -7\implies r\ge \sqrt{18}.\]

 The intersection of these inequalities is the circular region $T$ for which every circle in $T$ has a radius between $\sqrt{18}$ and $\sqrt{32}$, inclusive. The area of such a region is thus $\pi(32-18)=14\pi.$ The requested probability is therefore $\frac{14\pi}{128} = \frac{7\pi}{64},$ yielding $(m,n)=(7,64).$ We have $m+n=7+64=\boxed{\textbf{(B)}\ 71}.$

 -AbhiSood1234, countmath1

 



\textbf{Solution 2 (Calculus)}

Expressing the Area of Region \( B \)

 Region \( B \) consists of points where \( |x| + |y| \le 8 \)

 In each quadrant, this can be expressed by the following functions:

 First quadrant: \( y = 8 - x \)Second quadrant: \( y = 8 + x \)Third quadrant: \( y = -8 - x \)Fourth quadrant: \( y = -8 + x \)

 In the first quadrant, \( x \) ranges from 0 to 8, and \( y \) ranges from 0 to \( 8 - x \). Thus, the area in the first quadrant is:
 \[\text{Area of first quadrant} = \int_0^8 \int_0^{8 - x} \, dy \, dx\]
 \[= \int_0^8 [y]_{y=0}^{y=8-x} \, dx = \int_0^8 (8 - x) \, dx\]
 \[= \left[ 8x - \frac{x^2}{2} \right]_0^8 = 64 - 32 = 32\]The total area of region \( B \) is:
 \[\text{Area of } B = 4 \times 32 = 128\]

 Expressing the Area of Region \( T \)Region \( T \) is defined by the inequality \( (x^2 + y^2 - 25)^2 \le 49 \), which can be rewritten as:
 \[18 \le x^2 + y^2 \le 32\]

 To find the area, we switch to polar coordinates with \( x = r \cos \theta \) and \( y = r \sin \theta \), where \( x^2 + y^2 = r^2 \). Here, \( r \) ranges from \( \sqrt{18} \) to \( \sqrt{32} \), and \( \theta \) ranges from 0 to \( 2\pi \).

 The area of \( T \) can then be found by:
 \[\text{Area of } T = \int_0^{2\pi} \int_{\sqrt{18}}^{\sqrt{32}} r \, dr \, d\theta\]
 \[= \int_0^{2\pi} \left[ \frac{r^2}{2} \right]_{r=\sqrt{18}}^{r=\sqrt{32}} \, d\theta = \int_0^{2\pi} \left( \frac{32}{2} - \frac{18}{2} \right) \, d\theta\]
 \[= \int_0^{2\pi} 7 \, d\theta = 14\pi\]

 The probability \( P \) that a dart lands in region \( T \) is the area of \( T \) divided by the area of \( B \):
 \[P = \frac{\text{Area of } T}{\text{Area of } B} = \frac{14\pi}{128} = \frac{7\pi}{64}\]

 So the probability is of the form \( \frac{m}{n} \pi \), where \( m = 7 \) and \( n = 64 \), so \( m + n = 7 + 64 = 71 \).

 
 \[\boxed{\textbf{(B)}\ 71}\]

 By the way never pull out calculus on the amc10. Itis time consuming and easy to make mistakes on~[minor edits by hao11]

 



\textbf{Solution 3}


                    \begin{figure}[H]
                        \centering
                        \includegraphics[width=0.5\textwidth]{Solution/AMC_10_10B_2024_14_Solution_Figure_1.jpeg}
                    \end{figure}
                    


	\item (\textit{AMC 10 10B 2024}) A list of $9$ real numbers consists of $1$, $2.2$, $3.2$, $5.2$, $6.2$, and $7$, as well as $x$, $y$ , and $z$ with $x$ $\le$ $y$ $\le$ $z$. The range of the list is $7$, and the mean and the median are both positive integers. How many ordered triples ($x$, $y$, $z$) are possible?


 $\textbf{(A) } 1 \qquad \textbf{(B) } 2 \qquad \textbf{(C) } 3 \qquad \textbf{(D) } 4 \qquad \textbf{(E) } \text{infinitely many}$

 \textbf{Answer: }C



\textbf{Solution 1 (Main)}

$\textbf{First Case}$

 We start off by knowing that there must exist an ordered pair \( (0,y,z) \) and a big term 7 such that the range is satisfied and the median could also be satisfied. Because \( x \le y \le z \), we say \( x=0 \). Then, we get the sum 24.8. We need \( \frac{24.8 +y +z}{9} \in \mathbb{Z} \). This means that we need the nearest multiple of 9, which is 36, and we get \( y+z = 11.2 \). We need one of these to be the median, WLOG say \( y \). Then, \( y \in \{4,5\} \). So if \( y=4 \), we get \( z=7.2 \), and if \( y=5 \), we get 6.2. We see that if \( z=6.2 \), the range will still remain as 7, and therefore the one ordered triple \( (0,5,6.2) \) satisfies this.

 We know that the median can be \( y \), so we have \( y \in \{4,5\} \). We do casework on 4 and 5.

 $\textbf{Case 1}$

 Say \( y=4 \). Then, because we already know there exists an ordered triple where \( z \) is not the largest, there must exist at least one ordered triple where \( z \) is, so we say \( z-x=7 \). We also know that the mean must be divisible by 9. Quickly summing each number up we get 24.8. The next number divisible by 9 is 36. We add \( y \) to get 28.8. We then know \( x+z=36-28.8=7.2 \). The smallest values give \( x=0.1 \) and \( z=7.1 \). This satisfies our constraints.

 We don't want any errors, so we check the next sum, which is 45. We get then \( 7.2=x+z \) and \( z-x=7 \). This gives \( x=4.6 \) and \( z=11.6 \). This is a contradiction, as \( x \le y \le z \), and the values incrementally become too large, and therefore we only have one ordered triple for this case, which is \( (0.1,4,7.1) \)

 $\textbf{Case 2}$

 We now check \( y=5 \). This makes our sum 29.8. The nearest multiple of 9 is 36, again. This gives us \( 6.2 = x+z \) and \( z-x=7 \). Solving gives us -0.4 and 6.6. This however doesn't work, as 6.6 is not the greatest value.

 We decide to check the next multiple of 9, which is 45. This gives us \( 15.2 = x+z \) and \( z-x=7 \). This gives us \( x=4.1 \) and \( z=11.1 \). This also doesn't work, as \( x \) becomes incrementally larger, and therefore there are no ordered triples that satisfy this.

 $\textbf{Last Case}$

 We have one more case, and that is \(x\) is not the first number, but \( z \) is the last, meaning \( z-1=7 \), and \( z=8 \). This means that \( \frac{24.8+x+y+z}{9} \in \mathbb{Z} \). So, if \( z=8 \), we get 32.8. We again see that we need the nearest multiple of 9 which is 36, and we get \( 3.2 = x+y \). We see that \( x \) and \( y \) still need to be a integer median and 3.2 is too small, so 36 cannot work. We try 45, and see \( 12.2 = x+y \). After some trial and error, we see that \( x=6 \) and \( y=6.2 \), giving us the ordered triple \( (6,6.2,8) \). Continuing as before shows us that \( x \) again incrementally increases, and therefore we have only one ordered triple.

 We have $\boxed{\textbf{(C) }3}$ ordered triples. These are \( (0.1,4,7.1) \), \( (0,5,6.2) \), and \( (6,6.2,8) \).

 



\textbf{Solution 2 (Circular Arrangements)}

This is a more advanced solution, but it gives the answer really fast. Plus, it is always better to expand your knowledge.

 This solution is so overpowered that you don't even need to sum the decimals up. :-)

 Consider the arrangements like a circle. There are 3 terms; \( x \), \( y \), and \( z \), so we arrange them into a circle. We use the circular gap theorem to solve this, which is 
 \[C(nA) = \frac{n}{k} \binom{n-k-1}{k-1}\]

 We see that there are 3 terms. We need either 1 adjacent from the other, two adjacent from one, or three adjacent.

 Also keep in note that if the combinatorial part of \( C(nA) \) is invalid, it is just 0.

 We then have \( k \in \{1, 2, 3\} \).

 We then have 3 cases,

 $\textbf{Case 1:}$ $k \in \{1\}$

 We have \(  \frac{3}{1} \cdot \binom{3-1-1}{1-1} = \frac{3}{1} \cdot 1 = 3 \) ways to do this.

 $\textbf{Case 2:}$ $k \in \{2\}$

 We have \( \frac{3}{2} \cdot \binom{3-2-1}{2-1} \). The combinatorial part of this is \( \binom{0}{1} \), which is invalid and therefore 0. Therefore there are \( \frac{3}{2} \times 0 = 0 \) cases.

 $\textbf{Case 3:}$ $k \in \{3\}$

 We have \( \frac{3}{3} \cdot \binom{3-3-1}{3-1} \). The combinatorial part of this is \( \binom{-1}{2} \), and again, this is invalid, and therefore 0. There is then \( 1 \times 0 = 0 \) cases.

 Our answer is \( 3 + 0 + 0 =  \) $\boxed{\textbf{(C) }3}$

 \textbf{Disclaimer Solution 2} This theorem requires practice to use and when to use it. For example, if they said that this list is non circular or is strictly linear, then we may not solve it like this. However, these specifications have not been said in this problem.

 For disclaimer, don't use this formula everywhere. A clear example is that in a problem involving linear arrangements, this method will never work.

 Below are some practice problems that I highly recommend you try so you can understand this theorem and its applications better.

 Check them below.

 \href{https://artofproblemsolving.com/wiki/index.php/2014\_AMC\_10B\_Problems/Problem\_24}{2014 AMC 10B Q24}

 \href{https://artofproblemsolving.com/wiki/index.php/2017\_AMC\_10B\_Problems/Problem\_18?srsltid=AfmBOoorbWvdWI5mJYkf\_XsGNK8N7JEX7G2S4w13Uy\_lkqkFIXMbqJQc}{2017 AMC 10B Q18}

 \href{https://artofproblemsolving.com/wiki/index.php/1983\_AIME\_Problems/Problem\_7?srsltid=AfmBOoo2Rma8yq6kbV0holdmtaa2Zr2ezKPOD6cPyGa0Xs8GcMtwEf3J}{1983 AIME Q7}

 Also, it is helpful to note \( C_n(A) = \frac{n}{k} \binom{n-k-1}{k-1} \) = \( \binom{n}{2k} \).

 



\textbf{Solution 3 (AMC 10)}

The sum of the six existing numbers is $24.8$. For the mean to be an integer, the sum of the remaining three numbers $x,y,z$ and the remaining six must be divisible by $9$; thus $x+y+z$ must equal either $2.2,11.2,20.2$. However, one of $x,y,z$ must be the median since the middle four numbers in the original set are all non-integral. As a result, we can eliminate $2.2$ since the sum would be too small to allow for one of $x,y,z$ to be the median (notice that negative numbers are not permitted since the range would be greater than $7$). This leaves two cases.

 For the $x+y+z=20.2$ case, notice that all three must be less than or equal to $8$ due to the range restriction. However, also notice that $y$ and $z$ must be on the higher end of the range, meaning that the new median would have to fall between the new middle numbers $5.2$ and $6.2$. Thus $x=6$, and we can manipulate the numbers to make the range $7$ by making $z=8$ and thus $y=6.2$, providing one case in $(6,6.2,8)$.

 For the $x+y+z=11.2$ case, we note that the average of the three would fall near the median, signaling that one of the numbers would be the median, one would fall lesser than the median, and one would reach greater than the median. Thus we let the median be $y$. Since adding a number to each side of the set would not change the median, we know that the median must fall between the two middle numbers $3.2,5.2$. Then we find two triples $(0.1,4,7.1)$ and $(0,5,6.2)$, resulting in $\boxed{\textbf{(C) }3}$.

 



\textbf{Solution 4 (AMC 12)}

We can do casework on the range. We have four possible cases:

 $x<1$, $z<7$
$x>1$, $z<7$
$x<1$, $z>7$
$x>1$, $z>7$

 These encapsulate all possible values of $x$ and $z$ we can choose, so we're not leaving anything out. As we'll see, the problem will become simpler this way. Since we want integer mean, first note the sum of the values given to be $24.8$: when we add $x,y,z$, it must be a multiple of $9$ to yield an integer mean. Also remember that the median of an increasing list of $9$ numbers is the fifth number.

 $\textbf{Case 1:}$

 Since $x<1$, $z<7$, $x$ must be the minimum of this list of numbers and $7$ is the max: so the range is $7-x$. This must be equal to $7$, so $x=0$. Now we try to make the median an integer.With $x=0$, the fifth smallest number currently is $5.2$, which is clearly not an integer. But if we stick $y$ or $z$ in between $3.2$ and $5.2$, then the median will be an integer. So let one of $y,z$ be $4$: then the sum of all the numbers we have so far is $28.8$. To make this a multiple of $9$, the last number has to be $7.2$ to add up to $36$. However, we said that $z<7$, so this violates the bounds of this case. So then we set $y=5$, so our sum is $29.9$: then $z=36-29.9 = 6.2$. Here everything checks out, so for this case there is one solution: $(0,5,6.2)$.

 $\textbf{Case 2:}$

 Since $x>1$ and $z<7$, the minimum is $1$ and the maximum is $7$. Thus the range is always $6$, but we need the range to be $7$, so this case has no solutions.

 $\textbf{Case 3:}$

 We have $x<1, z>7$. Our minimum is $x$ and our maximum is $z$, so the range is $z-x = 7$. Since $x<1$, the median still falls on $5.2$, so $y$ can be either $4$ or $5$ to make the median an integer. For now, suppose $y=4$. Our sum is now $28.8+x+z$, which we set to the nearest multiple of $9:$ $36$. Thus $x+z = 7.2$. Now we have a system of equations: we add them to get $2z = 14.2 \implies z = 7.1 \implies x = 0.1$. Since the values of $z$ and $x$ satisfy our case, this is a possible solution.

 Now suppose $y=5$. Our sum is $29.8+x+z$, which must be equal to $36$, so $x+z = 6.2$. Adding the equations, we have $2z = 13.2 \implies z = 6.6$. But this violates our assumption that $z>7$, so this has no solutions.

 Therefore this case has one solution: $(0.1,4,7.1)$.

 $\textbf{Case 4:}$

 Since $x>1, z>7$, the max is $z$ and the min is $1$. Thus $z-1 = 7 \implies z = 8$. Now the median is between $5.2$ and $6.2$, so $x$ or $y$ must take the value $6$. Our sum is now $24.8+6+8 = 38.8$, and adding the last number has to equal $45$. Thus the last number is $6.2$, and this yields one solution: $(6,6.2,8)$.

 


 
Therefore we've found three solutions. These cases encompass all possible real values of $x$ and $z$, so we know we've left nothing out. The answer is $\boxed{\textbf{(C) }3}$




\textbf{Solution 5}

We can start doing casework on the numbers. Notice that the median can either be four, five or six because there has to be one or two numbers on one side. We can start bashing everything out, starting from letting $x=a, y=4, z=7+a$. Now plugging into the sequence, we get that $a, 1, 2.2, 3.2, 4, 5.2, 6.2, 7, 7+a$. Notice that the sum of the numbers not including $x, y, z$ is $24.8$. To make this mean an integer, $x+y+z$ must have a decimal $0.2$ when added up. Now, plugging in our values into the mean gets us $35.8+2a=9k$ for some integer $k$. Notice that $0\leq a\leq 1$, so the only value for $a$ is $0.1$, so we found our first ordered triplet, $(0.1, 4, 7.1)$. Let’s make the median $5$ now. Then we have to put $x$ and one side and $z$ on the other to “balance” the numbers, or make it so that there are 4 numbers on each side of $5$. Now, we can see that $x$ is $0$, and solving for $z$, we get that $z=6.2$. We get our second triple is $(0, 5, 6.2)$. Now coming to our third case scenario, we use $6$ as the median. Clearly it means that $x$ has to be the median, so that $y$ and $z$ can be on the same side as $x$. 
 \[1, 2.2, 3.2, 5.2, 6, 6.2, y, 7, z\] We will use this diagram. Notice that $z$ has to be $8$, and solving for $y$ yields $y=6.2$, so we found our third and last triplet, $(6, 6.2, 8)$, meaning that $\boxed{C}$ is the right answer.

 



\textbf{Solution 6 (Cheese)}

Answer choice E is too big and answer choices A,B are both too small, so get rid of A,B, and E. Then, after checking the answer choices respectively, you get 3 as the final answer. Therefore, the answer is answer choice $\boxed{\textbf{(C) }3}$.

 
Note: The answer to AMC problems is NEVER, EVER infinitely many. Also, the top part wasn't mine, it was somebody else. 

 



\textbf{Solution 7 (A (still interesting but slightly better) cheese)}

Ignore the fact that \( x \le y \le z \). Then we know that there should be some ordered pair \((0,y,z)\), \((x,0,z)\), and \((x,y,0)\), where we can permute \( x \), \( y \), or \( z \) such that they do satisfy \( x \le y \le z \). Our answer is then \(\boxed{\textbf{(C) }3}\).

 


	\item (\textit{AMC 10 10B 2024}) Jerry likes to play with numbers. One day, he wrote all the integers from $1$ to $2024$ on the whiteboard. Then he repeatedly chose four numbers on the whiteboard, erased them, and replaced them by either their sum or their product. (For example, Jerry's first step might have been to erase $1$, $2$, $3$, and $5$, and then write either $11$, their sum, or $30$, their product, on the whiteboard.) After repeatedly performing this operation, Jerry noticed that all the remaining numbers on the whiteboard were odd. What is the maximum possible number of integers on the whiteboard at that time?


 $\textbf{(A) } 1010 \qquad \textbf{(B) } 1011 \qquad \textbf{(C) } 1012 \qquad \textbf{(D) } 1013 \qquad \textbf{(E) } 1014$

 \textbf{Answer: }A



\textbf{Solution 1}

Consider the numbers as $1,0,1,0,...,1,0$. Note that the number of odd integers is monotonously decreasing.

 We need to get rid of $1012$ $0$'s, so we either add or multiply four $0$s together to get $1012\rightarrow 253 \rightarrow 63+1=64 \rightarrow 16 \rightarrow 4 \rightarrow 1.$

 To get rid of the final $0$, we need to consume three other $1$'s to result in one $1$. Thus the answer is $1012-2=\boxed{\textbf{(A) } 1010 }$,

 



\textbf{Solution 2}

(I will denote even numbers as $e$ and odd numbers as $\normalsize o$ for the purpose of this question, and $x$ as the answer we need)

 First, we observe that each time Jerry replaces four number with one (let's call this an "$operation$"), we lose exactly three numbers regardless of the type of $operation$ ($Additive$ or $Multiplicative$).

 $\Rightarrow x \equiv 2024 \equiv 2\pmod3$

 Already, we see that the only possible options are $\textbf{(A) } 1010$ or $\textbf{(D) } 1013$

 We need there to be no $e$'s on the board at the end, so we attempt to convert them into $\normalsize o$ while retaining as many existing $\normalsize o$'s as possible.

 Notice that converting 4 $e$'s into an $\normalsize o$ is not possible. We can apply an $Additive \, operation$ on 3 $e$'s and 1 $\normalsize o$ to get back 1 $\normalsize o$ ($e+e+e+\normalsize o=\normalsize o$) By doing so, we retain all the $\normalsize o$'s we use in an $operation$ (and don't generate any new ones). In the first $2024$ numbers, there are $1012$ $\normalsize o$'s. This means that $x \le 1012$, and the only option left is $\boxed{\textbf{(A) } 1010}$

 
Note: For the exact steps to arrive at 1010 odd numbers, perform

 $e+e+e+o=o$ (337 times)

 $e+o+o+o=o$ (1 time)

 



\textbf{Solution 3 (guessing strategy)}

Notice that with every step, Jerry removes $4$ numbers and adds $1$ number to the board. Thus, $3$ numbers are removed from the board in every step. Since $2024\equiv2\pmod{3}$, the answer choice must also be congruent to $2$ mod $3$. Only $1010$ and $1013$ among the answer choices satisfy this condition. This alone eliminates three choices, and one can guess from here for a 3 point expected value. 

 However, we start out with $1012$ even numbers and $1012$ odd numbers, and it is is not possible to get more odds than there were initially (every time we add or multiply four numbers, we can only produce one odd number, which cannot occur if only even numbers are used; therefore, there cannot be more odd numbers after an operation than before it). Thus, we choose $\boxed{\textbf{(A) } 1010.}$




\textbf{Solution 4 (OEEE)}

There are $1012$ odd numbers and $1012$ even numbers. In order to maximize the remaining numbers, then we must minimize the number of times the operation is performed.

 Since we want to be left with all odd numbers, then we must get rid of the even numbers as quick as possible, which implies $EEEE$. This doesn't give an odd product or sum, so our next best option is $OEEE$, which gives an odd sum. 

 The maximum number of times this operation can be performed with $OEEE$ is $\left\lfloor{\dfrac{1012}{3}}\right\rfloor=337$. This will eliminate $3\cdot337=1011$ even numbers and no odd numbers (because each time the operation is performed, one odd number is eliminated but another one is created). 

 Therefore after the operation is performed $337$ times with $OEEE$, there will be $1$ even number left and $1012$ odd numbers left. Performing this operation one more time with $EOOO$ will give $0$ even numbers and $\boxed{\text{(A) }1010}$ odd numbers. ~Tacos\_are\_yummy\_1

 



\textbf{Solution 5 (EEEE)}

There are $1012$ odd numbers and $1012$ even numbers. Our goal is to get rid of all the even numbers. By performing the operation with $EEEE$, we will remove $4$ even numbers and create one simultaneously.

 First, we can perform $EEEE$ a maximum of $\dfrac{1012}{4}=253$ times, to end up with $1012-4\cdot253+253=253$ even numbers. Performing $EEEE$ again $\left\lfloor{\dfrac{253}{4}}\right\rfloor=63$ times will result in $253-4\cdot63+63=64$ even numbers left. Repeating this process, eventually there will be $1$ even number left and still $1012$ odd numbers. Performing $OOOE$ gives the answer of $\boxed{\text{(A) }1010}$. ~Tacos\_are\_yummy\_1 ~ AVTV MINOR EDITS

 



\textbf{Solution 6 (Monovariants)}

Note that the number of odd integers is monotonically decreasing, and further observe that the number of odd integers is preserved after an operation where Jerry adds three evens and one odd together, producing an odd number. To maximise the number of odd integers, we will use this operation whenever possible.

 There are $1012$ even integers; to perform the specific procedure outlined above, we remove $3$ even integers each time. Therefore, we will arrive in a situation where there is $1$ even integers and $1012$ odd integers. To ensure that after the final operation, there will only be odd integers, we add the even integer and three odd integers, yielding one odd integer. There are now $1010$ integers, all of them odd, so we select $\fbox{ (A) 1010 }$.

 



\textbf{Solution 7}

There are $1012$ odd integers, so the choices D and E are not possible. If the numbers 1 and 3 were erased, they would not be able to be recreated, the answer is $1012$ - $2$ = $\fbox{ (A) 1010 }$.

 



\textbf{Solution 8}

Like Solution 4, we'll use $OEEE$, and they already explained why it maximizes the number of odd integers. Consider the set of integers between $1$ and $n$. Optimally, $n \equiv 0 \pmod{6}$ because $OEEE$ uses $3$ even integers (& the number of even integers is $\left \lfloor \frac{n}{2} \right\rfloor$ so $n$ is a multiple of $2$). In that case, the max number of odd integers remaining is $\frac{n}{2}$ where no odd integer is removed without being replaced.

 The closest multiple of $6$ to the number $2024$ is $2022$. A maximum of $1011$ odd integers can be left if the process is done for the integers $1$ through $2022$. Finally, there's $1$ even number left, so the operation $OOOE$ is performed, thus the answer is $1011 + 1 - 2 = \fbox{ (A) 1010 }$ since $2023$ is another odd integer not between $1$ and $2022$, but the operation $OOOE$ removes a net $2$ odd integers.

 



\textbf{Solution 9}

We can remove $3$ evens and $1$ odd every turn and do this $337$ times until we have $1$ even number remaining and $1012$ odd numbers because in every group of $4$ we lose $3$ evens, and one odd, but gain the odd back as the sum of the group will always be odd. With $1013$ integers left we remove a group of $4$ containing the one even and add back their odd some so $1013 - 4 + 1$ equals A, $\boxed{1010}$.

 


	\item (\textit{AMC 10 10B 2024}) In a race among 5 snails, there is at most one tie, but that tie can involve any number
of snails. For example, the result of the race might be that Dazzler is first; Abby,
Cyrus, and Elroy are tied for second; and Bruna is fifth. How many different results
of the race are possible?


 $\textbf{(A) } 180 \qquad\textbf{(B) } 361 \qquad\textbf{(C) } 420 \qquad\textbf{(D) } 431 \qquad\textbf{(E) } 720$

 \textbf{Answer: }D



\textbf{Solution 1}

Intuitive solution:We perform casework based on how many snails tie. Let's say we're dealing with the following snails: $A,B,C,D,E$.

 $5$ snails tied: All $5$ snails tied for $1$st place, so only $1$ way.

 $4$ snails tied: $A,B,C,D$ all tied, and $E$ either got $1$st or last. ${5}\choose{1}$ ways to choose who isn't involved in the tie and $2$ ways to choose if that snail gets first or last, so $10$ ways.

 $3$ snails tied: We have $ABC, D, E$. There are $3! = 6$ ways to determine the ranking of the $3$ groups. There are $5\choose2$ ways to determine the two snails not involved in the tie. So $6 \cdot 10 = 60$ ways.

 $2$ snails tied: We have $AB, C, D, E$. There are $4! = 24$ ways to determine the ranking of the $4$ groups. There are $5\choose{3}$ ways to determine the three snails not involved in the tie. So $24 \cdot 10 = 240$ ways.

 $1$ snail tied: This is basically just every snail for a place, so $5! = 120$ ways.

 The answer is $1+10+60+240+120 = \boxed{\text{(D) }431}$.

 



\textbf{Solution 2}

Split the problem into cases. A tie of $n$ snails has $\dbinom{5}{n}$ ways to choose the snails that are tied, $6-n$ ways to choose which place they tie for, and $(5-n)!$ to place the remaining snails.

 1. No tie $\implies5!=120$

 2. Tie of 2 snails $\implies\dbinom{5}{2}\cdot4\cdot3!=240$

 3. Tie of 3 snails $\implies\dbinom{5}{3}\cdot3\cdot2!=60$

 4. Tie of 4 snails $\implies\dbinom{5}{4}\cdot2=10$

 5. Tie of all 5 snails $\implies1$

 The answer is $120+240+60+10+1=\boxed{\text{(D) }431}$ ~Tacos\_are\_yummy\_1

 



\textbf{Solution 3 (Answer Choice Reasoning)}

For the case of a 5-way tie, we have $\dbinom{5}{5}=1$ cases. We can assume this leads to an answer that ends in 1, leaving only $B$ and $D$. By just accounting for the cases for a 1-way tie ($5!=120$) and 2-way tie ($\dbinom{5}{2}\cdot4\cdot3!=240$), it immediately shows that $B$ does not work. Therefore, the answer is $\boxed{\text{(D) }431}$~shreyan.chethan, cleaned up by cweu001, vintage\_god

 Notice that with 2 remaining choices, through guessing, the expected value of your points is $\frac{6}{2}=3$, and not answering gives 1.5 points. Therefore, you gain an expected value of $1.5$ points by answering. It would be best to always guess whenever you can get rid of $2$ or more answers. If you can only eliminate $1$ answer, not guessing will give you $1.5$ points, but guessing also gives you an expected value of $1.5$ points. I usually leave it blank if I can only get rid of $1$ answer choice.~shreyan.chethan, edited by BenjaminDong01, added by PerseverePlayer , and added by Aarav22

 Sidenote : Usually, you should aim to look at the closer option choices than the extreme ones.

 



\textbf{Solution 4 (Rigorous Solution 1; same logic)}

Call the snails $A,B,C,D,E$ in said order apart of the set $S$, and define $S_{T}$ to be the set of ties regarding $T$ snails. Define a \textbf{maxima tie} to be a set $\mathcal{M}$ that retains a tie between $k$ members for $2 \le k \le 5$. Consider casework on all \textbf{maxima ties}.

 If no \textbf{maxima ties} exist, then the problem degenerates into how many word rearrangements exist for $ABCDE$, which is $5! = 120$.

 If $k=2$ \textbf{maxima tie} exist, then note that if the tie is held between $A,B$, it is the same as a tie holding for $B,A$, meaning we have $\frac{5 \cdot 4}{2} = 10$ ways to choose our \textbf{maxima tie}. Then, the word rearrangement for $\mathcal{M}CDE$ is $4! = 24$, yielding 240 cases.

 If $k=3$ \textbf{maxima tie} exist, then note that if a tie is held between $A,B,C$, it is the same for the other 5 iterations. Out of one iteration, 5 other copies exist that yield the same result, meaning our cases for the \textbf{maxima tie} are $\frac{5 \cdot 4 \cdot 3}{6} = 10$. Then, the word rearrangement for $\mathcal{M}DE$ is $3! = 6$, yielding 60 cases.

 If $k=4$ \textbf{maxima tie} exist, then every partition of $A,B,C,D$ yields 23 other copies, meaning there are 5 acknowledgeable cases. Then the word rearrangement for $\mathcal{M}E$ is $2! = 2$, yielding 10 cases.

 If $k=5$ \textbf{maxima tie} exist, then every partition of $A,B,C,D,E$ yields 119 other copies, meaning there is only one acknowledgeable case. Then, the word rearrangement for $\mathcal{M}$ is $1! = 1$ yielding 1 case.

 The desired answer is then

 
 \[120 + \sum_{T = k \neq {1} \in S_{T}} \textbf{maxima tie plausible cases}\]
 \[120 + 240 + 60 + 10 + 1 = \boxed{431}\]



	\item (\textit{AMC 10 10B 2024}) How many different remainders can result when the $100$th power of an integer is
divided by $125$?


 $\textbf{(A) } 1 \qquad\textbf{(B) } 2 \qquad\textbf{(C) } 5 \qquad\textbf{(D) } 25 \qquad\textbf{(E) } 125$

 \textbf{Answer: }B



\textbf{Solution 1}

First note that the \href{https://artofproblemsolving.com/wiki/index.php?title=Euler\%27s\_totient\_function}{Euler's totient function} of $125$ is $100$. We can set up two cases, which depend on whether a number is relatively prime to $125.$

 If $n$ is relatively prime to $125$, then $n^{100} \equiv 1 \pmod{125}$ because of \href{https://artofproblemsolving.com/wiki/index.php?title=Euler\%27s\_Totient\_Theorem}{Euler's Totient Theorem}. 

 If $n$ is not relatively prime to $125$, it must have a factor of $5$. Express $n$ as $5m$, where $m$ is some integer. Then $n^{100} \equiv (5m)^{100} \equiv 5^{100}\cdot m^{100} \equiv 125 \cdot 5^{97} \cdot m^{100} \equiv 0 \pmod{125}$.

 Therefore, $n^{100}$ can only be congruent to $0$ or $1 \pmod{125}$. Our answer is $\boxed{2}$.

 \textbf{Sidenote: Proof of Euler's Theorem} Euler's Theorem states that $a^{\phi(n)}\equiv 1\pmod{n}$, where $\phi(n)$ is Euler's totient function, which counts the number of positive integers $x$ such that $1\le x \le n$ and $\gcd(x,n)=1$. Note that this only works for positive integer $a$ such that $\gcd(a,n)=1.$

 Consider the set of distinct positive integers $S=\{x_1,x_2,\ldots x_{\phi(n)}\}$ such that $1\le x_i\le n$ and $\gcd(x_i,n)=1.$ Multiply all terms by $a$ to create $aS=\{ax_1,ax_2,\ldots ax_{\phi(n)}\}.$ First, we prove that $\gcd(ax_i,n)=1$ for any $x_i$ using contradiction. Let $p$ be a prime such that $p \mid ax_i$ and $p \mid n.$ This means that either $p \mid a$ or $p \mid x_i,$ which is not possible since both $\gcd(a,n)$ and $\gcd(x_i,n)=1.$  Next we show that no two elements in set $aS$ are congruent modulo $n.$ Let $x_i$ and $x_j$ exist such that $ax_i\equiv ax_j \pmod{n}.$ Multiplying both sides by $a^{-1}$(which exists) gives $x_i\equiv x_j\pmod{n}.$ This means that if $x_i \not\equiv x_j \pmod{n},$ then $ax_i \not\equiv ax_j \pmod{n}.$ Now we multiply all terms in each set, and note that each residue appears exactly once in both sets, so we have: 
 \[a^{\phi(n)}\cdot \prod_{i=1}^{\phi(n)} x_i \equiv \prod_{i=1}^{\phi(n)} x_i \pmod{n}.\] Because each $x_i$ has the property that $\gcd(x_i,n)=1,$ we must have that $\gcd(\prod_{i=1}^{\phi(n)} x_i,n)=1.$ Then multiplying both sides by the inverse of product thereof gives $a^{\phi(n)}\equiv 1\pmod{n}.$ $\blacksquare$




\textbf{Solution 2 (No Totient)}

Note that

 
 \[(5+n)^{100} = {100 \choose 0} 5^{100} + {100 \choose 1} 5^{99} n + {100 \choose 2} 5^{98} n^2 + \cdots + {100 \choose 97} 5^3 n^{97} + {100 \choose 98} 5^2 n^{98} + {100 \choose 99} 5 n^{99} + {100 \choose 100} n^{100}.\]

 Taking this mod $125$, we can ignore most of the terms except the for the last $3$:

 
 \[{100 \choose 98} 5^2 n^{98} + {100 \choose 99} 5 n^{99} + {100 \choose 100} n^{100} \equiv 4950 \cdot 5^2 n^{98} + 100 \cdot 5 n^{99} + n^{100} \equiv n^{100} \pmod {125},\]

 so $(5+n)^{100} \equiv n^{100}$. Substituting $-n$ for $n$, we get $(5-n)^{100} \equiv n^{100}$. Therefore, the remainders when divided by $125$ repeat every $5$ integers, so we only need to check the $100$th powers of $0, 1, 2, 3, 4$. But we have that $(5-1)^{100} \equiv 1^{100}$ and $(5-2)^{100} \equiv 2^{100}$, so we really only need to check $0, 1, 2$. We know that $0, 1$ produce different remainders, so the answer to the problem is either $2$ or $3$. But $3$ is not an answer choice, so the answer is $\boxed{\textbf{(B) } 2}$.

 



\textbf{Solution 3 (Binomial Theorem)}


                    \begin{figure}[H]
                        \centering
                        \includegraphics[width=0.5\textwidth]{Solution/AMC_10_10B_2024_18_Solution_Figure_1.jpeg}
                    \end{figure}
                    



\textbf{Solution 4}

Firstly, I interpreted the problem algebraically. In other words, let ${x}$ be an integer. Since we're raising it to the ${100}$th power and ${100}$ is even, we don't have to worry about negative values of ${x}$. Now, we break ${125}$ into three factors of ${5}$. I first considered this when we're dividing by ${5}$. Okay so we can write the problem as below:

 $x^{100}\equiv r \bmod{5}$. Now, in the world of mod(5), everything is much easier to figure out. Take $x\equiv 0 \bmod{5}$. Now, obviously $x^{100}$ will bring up a remainder of ${0}$ if ${x}$ itself is a multiple of 5. Now, $x\equiv 1 \bmod{5}$. In this case, notice how $1^{100}$ is congruent to 1 in the mod(5) world. So in this case, we would have a remainder of 1. Considering $x\equiv 2 \bmod{5}$, it's a little more tricky. Notice $2^{100}$ is equal to $1024^{10}$ and ${1024}$ is 1 less than a multiple of 5. So ${1024}$ will be congruent to  $-1 \bmod{5}$ and raising it to the power of 10 will bring up a remainder of 1 yet again. Similarly, is $x\equiv 3 \bmod{5}$ or $x\equiv 4 \bmod{5}$, you would notice both these cases will bring up a remainder of 1 again. So in the mod(5) world, we seem to only have remainders of ${0}$ and ${1}$. Interesting!Therefore, since ${125}$ is just ${3}$ factors of ${5}$, we may say that in mod(125), we would have the same exact remainders namely ${0}$ and ${1}$. Therefore, we would only have ${2}$ possible remainders. This brings an answer of $\boxed{\textbf{(B)} \; 2}$.

 


	\item (\textit{AMC 10 10B 2024}) In the following table, each question mark is to be replaced by "Possible" or "Not
Possible" to indicate whether a nonvertical line with the given slope can contain the
given number of lattice points (points both of whose coordinates are integers). How
many of the 12 entries will be "Possible"?


 
                    \begin{figure}[H]
                        \centering
                        \includegraphics[width=0.5\textwidth]{Problem/AMC_10_10B_2024_19_Problem_Figure_1.png}
                    \end{figure}
                    
$\textbf{(A) } 4 \qquad     \textbf{(B) } 5 \qquad \textbf{(C) } 6    \qquad \textbf{(D) } 7      \qquad \textbf{(E) } 9$

 \textbf{Answer: }C



\textbf{Solution 1}

Let's look at zero slope first. All lines of such form will be expressed in the form $y=k$, where $k$ is some real number. If $k$ is an integer, the line passes through infinitely many lattice points. One such example is $y=1$. If $k$ is not an integer, the line passes through $0$ lattice points. One such example is $y=1.1$. So we have $2$ cases.

 
Let's now look at the second case. The line has slope $\frac{p}{q}$, where $p$ and $q$ are relatively prime integers. The line has the form $y = \frac{p}{q}x + m$. If the line passes through lattice point $(a,b)$, then the line must also pass through the lattice point $(a+kq, b+kp)$, where $a,b,k$ are all integers. Therefore, the line can pass through infinitely many lattice points but it cannot pass through exactly $1$ or $2$. The line can pass through $0$ lattice points, such as $y=x+0.5$. This contributes $2$ more cases.

 
If the line has an irrational slope, it can never pass through more than $2$ lattice points. We prove this using contradiction. Let's say the line passes through lattice points $(a,b)$ and $(c,d)$. Then the line has slope $\frac{d-b}{c-a}$, which is rational. However, the slope of the line is irrational. Therefore, the line can pass through at most $1$ lattice point. One example of this is $y=\sqrt{2}x$. This line contributes $2$ final cases.

 Our answer is therefore $\boxed{6}$.

 This shows examples of the six lines that are valid. - \href{https://www.desmos.com/calculator/ypt7esdv2g}{https://www.desmos.com/calculator/ypt7esdv2g}

 



\textbf{Solution 2}

For a 0 slope: It can have 0 points (e.g. $y=1.5$), but it can also have infinite (e.g. y=1), which checks the two or more box.

 For a nonzero slope: It can have 0 lattice points (e.g. $y=\frac23 x + \frac12$), or infinite (eg $y=2x$).

 For an irrational slope: It can have 0 (e.g. $y=\sqrt3 x+\frac12$) or exactly 1 (e.g. $y=\sqrt3 x,$ which has an intersection at 0,0)

 Hence, there are a total of 6 points.

 



\textbf{Solution 3}

Zero Slope

 A line of zero slope can intersect no lattice points (If the y-intercept isn't an integer)

 It cannot intersect only one lattice point, same goes for two lattice points because if it intersects one, it has to intersect all lattice points with that y value.

 It can intersect more than two lattice points (infinite)

 Rational Slope $\neq0$

 It can intersect no lattice points ($y=x+0.5$)

 It cannot intersect only one or two lattice points (because then it would intersect another one once the slope "repeats" again)

 It can intersect more than two lattice points (infinite)

 Irrational Slope

 It can intersect no lattice points (y-intercept isn't an integer)

 It can intersect one lattice point (y-intercept is an integer)

 It cannot intersect two or more lattice points no matter what.

 Therefore, the answer is $\boxed{\text{(C) }6}$. ~Tacos\_are\_yummy\_1

 


	\item (\textit{AMC 10 10B 2024}) Three different pairs of shoes are placed in a row so that no left shoe is next to a
right shoe from a different pair. In how many ways can these six shoes be lined up?


 $\textbf{(A) } 60 \qquad\textbf{(B) } 72 \qquad\textbf{(C) } 90 \qquad\textbf{(D) } 108 \qquad\textbf{(E) } 120$

 \textbf{Answer: }A



\textbf{Solution 1}

Let $A_R, A_L, B_R, B_L, C_R, C_L$ denote the shoes.

 
There are $6$ ways to choose the first shoe. WLOG, assume it is $A_R$. We have $A_R,$ \_\_, \_\_, \_\_, \_\_, \_\_.

 
$~~~~~$ Case $1$: The next shoe in line is $A_L$. We have $A_R, A_L,$ \_\_, \_\_, \_\_, \_\_. Now, the next shoe in line must be either $B_L$ or $C_L$. There are $2$ ways to choose which one, but assume WLOG that it is $B_L$. We have $A_R, A_L, B_L,$ \_\_, \_\_, \_\_.

 
$~~~~~ ~~~~~$ Subcase $1$: The next shoe in line is $B_R$. We have $A_R, A_L, B_L, B_R,$ \_\_, \_\_. The only way to finish is $A_R, A_L, B_L, B_R, C_R, C_L$.

 
$~~~~~ ~~~~~$ Subcase $2$: The next shoe in line is $C_L$. We have $A_R, A_L, B_L, C_L,$ \_\_, \_\_. The only way to finish is $A_R, A_L, B_L, C_L, C_R, B_R$.

 
$~~~~~$ In total, this case has $(6)(2)(1 + 1) = 24$ orderings.

 
$~~~~~$ Case $2$: The next shoe in line is either $B_R$ or $C_R$. There are $2$ ways to choose which one, but assume WLOG that it is $B_R$. We have $A_R, B_R,$ \_\_, \_\_, \_\_, \_\_.

 
$~~~~~ ~~~~~$ Subcase $1$: The next shoe is $B_L$. We have $A_R, B_R, B_L,$ \_\_, \_\_, \_\_.

 
$~~~~~ ~~~~~ ~~~~~$ Sub-subcase $1$: The next shoe in line is $A_L$. We have $A_R, B_R, B_L, A_L,$ \_\_, \_\_. The only way to finish is $A_R, B_R, B_L, A_L, C_L, C_R$.

 
$~~~~~ ~~~~~ ~~~~~$ Sub-subcase $2$: The next shoe in line is $C_L$. We have $A_R, B_R, B_L, C_L,$ \_\_, \_\_. The remaining shoes are $C_R$ and $A_L$, but these shoes cannot be next to each other, so this sub-subcase is impossible.

 
$~~~~~ ~~~~~$ Subcase $2$: The next shoe is $C_R$. We have $A_R, B_R, C_R,$ \_\_, \_\_, \_\_. The next shoe in line must be $C_L$, so we have $A_R, B_R, C_R, C_L,$ \_\_, \_\_. There are $2$ ways to finish, which are $A_R, B_R, C_R, C_L, A_L, B_L$ and $A_R, B_R, C_R, C_L, B_L, A_L$.

 
$~~~~~$ In total, this case has $(6)(2)(1 + 2) = 36$ orderings.

 
Our final answer is $24 + 36 = \boxed{\textbf{(A) } 60}$




\textbf{Solution 2 (Tree Diagram)}

Label the pairs \( (L_1, R_1) \), \( (L_2, R_2) \), and \( (L_3, R_3) \).

 We now create a tree diagram. \href{https://artofproblemsolving.com/wiki/index.php?title=WLOG}{WLOG}, we say our first tree diagram starts with \( L_1 \).

 
                    \begin{figure}[H]
                        \centering
                        \includegraphics[width=0.5\textwidth]{Solution/AMC_10_10B_2024_20_Solution_Figure_1.png}
                    \end{figure}
                    This tree has 10 branches.

 Notice how we can interchange the first shoe with any \( L_1 \), \( L_2 \), \( L_3 \), \( R_1 \), \( R_2 \), \( R_3 \), and still get the same tree, but just interchanged shoes. Therefore, we just multiply the tree's branches by 6, to get \( 6 \cdot 10 =  \) $\boxed{\textbf{(A) } 60}$.

 



\textbf{Solution 3 (Fakesolve, Refer to Solution 4)}

We have two main cases: 

 \textbf{Case 1:} $LLLRRR$ or $RRRLLL$ 

 There are $3! = 6$ ways to assign the shoes to the ordering. For now we'll focus on the $LLLRRR$ case and multiply by $2$ since they're symmetrical. WLOG say we have $A_LB_LC_LC_RA_RB_R$. We can flip $A_L$ and $B_L$ among themselves and $A_R$ and $B_R$ among themselves, so we get $6 \cdot 2 \cdot 2 = 24.$ Multiplying by $2$ for symmetry gives us $48$ ways for this case. 

 


 \textbf{Case 2:} $RRLLLR$ or $RLLLRR$ 

 Again, we'll focus on the $RRLLLR$ case and multiply by $2$ since they're symmetrical. WLOG say we have $B_RA_RA_LB_LC_LC_R$. We can assign the shoes to the ordering in $3! = 6$ ways. $B_R$ can also be moved over to be next to $C_R$. Then we would have $A_RA_LB_LC_LC_RB_R$, which is also valid. So this case gives us $6 \cdot 2 = 12$ ways. 

 The answer is $48 + 12 = \boxed{60}.$ 

 Note: Case 1 should give you 24 and case 2 should give you 36

 



\textbf{Solution 4(focus on restrictions)}

Notice that you cannot have $LRL$ or $RLR$ in a row, since you are guaranteed an $R$ and an $L$ from a different pair. This means you can either have three $L$'s in a row, three $R$'s in a row, or you have two $R$'s between two $L$'s and two $L$'s between two $R$'s. 

 Below are the cases(note that once an $L$ is fixed the $R$ adjacent to it is also fixed due to the constraint):\begin{align*}LLLRRR \Rightarrow 3!\cdot 2!=12\\RLLLRR \Rightarrow 3!=6\\RRLLLR \Rightarrow 3!=6\\RRRLLL \Rightarrow 3!\cdot 2!=12\\LRRRLL \Rightarrow 3!=6\\LLRRRL \Rightarrow 3!=6\\LRRLLR \Rightarrow 3!=6\\RLLRRL \Rightarrow 3!=6\\\end{align*}


 We have $2\cdot 12+6\cdot 6=\boxed{\textbf{(A) }60}.$




\textbf{Solution 5(for some reason this works)}

Let the pairs be $L_1 R_1, L_2 R_2, L_3 R_3$. The only ones that don't work are

 $L_1 R_2, L_1 R_3, L_2 R_1, L_2 R_3, L_3 R_2, L_3 R_1$ and their swaps so a total of $12$ pairs that can't be in the string of $6$. We have $6! = 720$ total ways and since there are $12$ ways it won't work we simply do

 $\frac{720}{12} = \boxed{60}$




\textbf{Solution 6}

Let $A_1, A_2, A_3, B_1, B_2, B_3$ be the shoes. Initially, let's assume that it doesn't matter whether shoes of type $A$ are left/right or if shoes of type $B$ are left/right. We can have four possible arrangements: $AAABBB$, $AABBBA$, $ABBAAB$, and $ABBBAA$. Note that whenever we have a shoe of type $A$ and a shoe of type $B$ next to each other, those shoes must be of the same pair.

 In the first case ($AAABBB$), only one pair is forced to be in the middle. We can choose this pair in $3$ ways. There are $2! \cdot 2! = 4$ ways to choose the different pairs for the other four shoes in this case, and therefore $3 \cdot 4 = 12$ possibilities for this case.

 In all of the other cases, we either have $2$ or $3$ forced pairs $AB$. In all of these cases, there are $3 \cdot 2$ or $3 \cdot 2 \cdot 1$ ways to choose the different pairs that each shoe belongs to, so $6$ cases for each pair.

 In total, there are $12+6+6+6=30$ cases ignoring if $A$ or $B$ is left/right. If we now consider these different cases, we will have twice the number of combinations, or $30 \cdot 2 = \boxed{60}$ combinations.

 



\textbf{Solution 7}

Same thing for solution 1 but you use casework on whether all L's touch or not.Case 1: All L's touch24+12=36 possibilitiesCase 2: Only 2 L's touch24 possibilities36+24=$60$. 

 


	\item (\textit{AMC 10 10B 2024}) Two straight pipes (circular cylinders), with radii $1$ and $\frac{1}{4}$, lie parallel and in contact on a flat floor. The figure below shows a head-on view. What is the sum of the possible radii of a third parallel pipe lying on the same floor and in contact with both?


 
                    \begin{figure}[H]
                        \centering
                        \includegraphics[width=0.5\textwidth]{Problem/AMC_10_10B_2024_21_Problem_Figure_1.png}
                    \end{figure}
                    
$\textbf{(A)}~\frac{1}{9} \qquad\textbf{(B)}~1 \qquad\textbf{(C)}~\frac{10}{9} \qquad\textbf{(D)}~\frac{11}{9} \qquad\textbf{(E)}~\frac{19}{9}$

 \textbf{Answer: }C



\textbf{Solution 1}

Notice that the sum of radii of two circles tangent to each other will equal to the distance from center to center. Set the center of the big circle be at $(0,1).$ Since both circles are tangent to a line (in this case, $y=0$), the y-coordinates of the centers are just its radius. 

 Hence, the center of the smaller circle is at $\left(x_2, \frac14\right)$. From the the sum of radii and distance formula, we have:
 \[1+\frac14 = \sqrt{x_2^2 + \left(\frac34\right)^2} \Rightarrow x_2 = 1.\]

 So, the coordinates of the center of the smaller circle are $(1, \frac14).$ Now, let $(x_3,r_3)$ be the coordinates of the new circle. Then, from the fact that sum of radii of this circle and the circle with radius $1$ is equal to the distance from the two centers, you have:
 \[\sqrt{(x_3-0)^2+(r_3-1)^2} = 1+r_3.\]Similarly, from the fact that the sum of radii of this circle and the circle with radius $\frac14$, you have:
 \[\sqrt{(x_3-1)^2+\left(r_3-\frac14\right)^2}= \frac14 + r_3.\]Squaring the first equation, you have:
 \[x_3^2+r_3^2-2r_3+1=1+2r_3+r_3^2 \Rightarrow 4r_3=x_3^2 \Rightarrow x_3=2\sqrt{r_3}.\]Squaring the second equation, you have:
 \[x_3^2-2x_3+1+r_3^2-\frac{r_3}{2}+\frac{1}{16}=\frac{1}{16}+\frac{r_3}{2}+r_3^2 \Rightarrow x_3^2-2x_3+1=r_3\]Plugging in from the first equation we have
 \[r_3-1=x_3^2-2x_3=4r_3-4\sqrt{r_3} \Rightarrow 3r_3-4\sqrt{r_3}+1=0 \Rightarrow (3\sqrt{r_3}-1)(\sqrt{r_3}-1)=0 \Rightarrow r_3=1, \frac19.\]Summing these two yields $\boxed{\frac{10}{9}}.$

 \textbf{Diagram} 



\textbf{Solution 2 (Clever Manipulation)}

In general, let the left and right outer circles and the center circle have radii $r_1,r_2,r_3$. When three circles are tangent as described in the problem, we can deduce $\sqrt{r_3}=\frac{\sqrt{r_1r_2}}{\sqrt{r_1}+\sqrt{r_2}}$ by Pythagorean theorem. 

 Setting $r_1=1, r_2=\frac14$ we have $r_3=\frac19$, and setting $r_1=1,r_3=\frac14$ we have $r_2=1$. Thus our answer is $\boxed{\textbf{(C)}}$. 

 



\textbf{Solution 3 (Descartes’s Theorem)}

Descartes’s Theorem states that for curvatures $k_1,k_2,k_3,k_4$ we have

 
 \[(k_1+k_2+k_3+k_4)^2=2(k_1^2+k_2^2+k_3^2+k_4^2)\]

 with a curvature being the reciprocal of the radius of a circle, being positive if it is externally tangent to other circles, negative if it is internally tangent to other circles, and zero if we consider a line as a degenerate circle. Here our curvatures are $k_1=1,k_2=4,k_3=0$, and we wish to find $k_4$. Plugging these values into our formula yields:

 
 \[(1+4+0+k_4)^2=2(1^2+4^2+0^2+k_4^2)\]

 \[(5+k_4)^2=2(17+k_4^2)\]

 \[k_4^2+10k_4+25=2k_4^2+34\]

 \[k_4^2-10k_4+9=0\]

 \[k_4=1,9\]

 The curvature and the radius are reciprocals, so our radii must be $1$ and $\frac{1}{9}$, and their sum is $\boxed{\textbf{(C) }\frac{10}{9}}$.

 



\textbf{Solution 4 (Advanced Guesswork)}

Looking at the diagram, we see that there must be a circle with a radius smaller than $\frac14$, and there must be a circle with a radius close to $1$. Looking at the answers, we can assume that the answer choice is above $1$, eliminating the first two answers. We can also assume that some students will only solve for the smaller circle, while others may only solve for the larger circle. It would be reasonably safe to say that one of the answers must be the radius of the smaller circle, while another must be the radius of the bigger circle. Looking at the answers, we see that  $\frac19$ and $1$ are reasonably close to the radii of the missing circles in the diagram, so we add them up to get $\boxed{\textbf{(C) }\frac{10}{9}}$.

 



\textbf{Solution 5 (Essentially Solution 3)}

For a problem like this with three externally tangent circles that are all tangent to a line, we can use a specific form ofDescartes's Circle Theorem. Let the curvatures of the three circles be $r_1, r_2, r_3$. The curvature of a circle is justthe reciprocal of the radius. Then we get, 
 \[(r_1 + r_2 + r_3)^2 = 2(r_1^2 + r_2^2 + r_3^2)\] Plugging in the information we already have, we get:
 \[(1 + 4 + r_3)^2 = 2(1^2 + 4^2 + r_3^2)\]
 \[(5 + r_3)^2 = 2(17 + r_3^2)\]
 \[25 + 10r_3 + r_3^2 = 34 + 2r_3^2\]
 \[r_3^2 - 10r_3 - 9 = 0\]
 \[(r_3 - 1)(r_3 - 9) = 0\]So, we get our two curvatures: 1 and 9. The radii equivalent are $1$ and $\frac{1}{9}$. Adding these up, we get $\boxed{\textbf{(C) }\frac{10}{9}}$.

 



\textbf{Solution 6 (Metasolving)}

By looking at the answer choices, we realize that the answers are all sums of each other, and thus think "The AMC problem committee made the answers $\frac{1}{9}$ and $1$ so the people who only got one of the answers will pick either of the first two answers and those who counted one of the answers twice would pick either of the last two answers", and realize that the answers are most likely $1$ and $\frac{1}{9}$. Confirming with a ruler, we get $\frac{1}{9}$ + $1$ = $\boxed{\textbf{(C) }\frac{10}{9}}$, which is the solution. 

 $\sim$Glowworm

 Minor issue with this is that 10/9 is also an answer choice, and 1/9 + 10/9 = 11/9 which is also an answer choice, so it comes down to a bit of luck.

 



\textbf{Solution 7(practically similar to the other solutions)}

We can immediately see that we can put a circle with radius $1$ outside of the two circles already drawn that are tangent to the circles. If you brought a ruler with you, you can confirm this. We now find the radius of the circle between the circles of radius $1$ and $\frac{1}{4}$.. By drawing out the diagram, a geometry hack is to draw lines connecting the centers of the circles. Using the fact that tangent lines to the radius are perpendicular, we can construct a trapezoid with heights of lengths $1$ and $\frac{1}{4}$, a slant base of length $1 + \frac{1}{4} = \frac{5}{4}$. To find the base of the trapezoid, we can construct a right triangle with the slant base as the hypotenuse and a leg of length $1 - \frac{1}{4} = \frac{3}{4}$. Therefore, we see we have a $3-4-5$ right triangle and the length of the base of the trapezoid is conveniently just $1$. Let the radius of the circle between the two given circles $r$. By drawing the center of this circle and connecting the radius to the radii of the other two circles, we can construct two more trapezoids similar to the initially big one we constructed: one with perpendicular heights of lengths $1$ and $r$ and a slant base of $1 + r$, and another with perpendicular heights of lengths $\frac{1}{4}$ and $r$ and a slant base of length $r + \frac{1}{4}$. To find the bases of these trapezoids,(which we know sums to $1$ as we found earlier), we can use a similar trick by constructing right triangles. We can construct two right triangles: one with hypotenuse $1 + r$ and a leg of length $1 - r$, and another with hypotenuse $r + \frac{1}{4}$ and a leg of length $\frac{1}{4} - r$. We can use the Pythagorean Theorem to find the other legs of these right triangles and we can set the sum to equal $1$. Conveniently, the leg lengths turns out of be $2\sqrt{r}$ and $\sqrt{r}$, respectively, and therefore the sum will just be $3\sqrt{r}$. Since this is equal to $1$, we can solve $3\sqrt{r} = 1$ to get the radius to be equal to $\frac{1}{9}$. Now, we add $1$ and $\frac{1}{9}$ to get $\frac{10}{9}$ and our answer is just $\boxed{\textbf{(C) }\frac{10}{9}}$.

 
                    \begin{figure}[H]
                        \centering
                        \includegraphics[width=0.5\textwidth]{Solution/AMC_10_10B_2024_21_Solution_Figure_1.png}
                    \end{figure}
                    

 credit to Elephant200 for the diagram

 


	\item (\textit{AMC 10 10B 2024}) A group of $16$ people will be partitioned into $4$ indistinguishable $4$-person
committees. Each committee will have one chairperson and one secretary. The
number of different ways to make these assignments can be written
as ${3^r}M$, where $r$ and $M$ are positive integers and $M$ is not divisible by $3$. What is $r$?


 $\textbf{(A) } 5 \qquad\textbf{(B) } 6 \qquad\textbf{(C) } 7 \qquad\textbf{(D) } 8 \qquad\textbf{(E) } 9$

 \textbf{Answer: }A



\textbf{Solution 1}

There are ${16 \choose 4}$ ways to choose the first committee, ${12 \choose 4}$ ways to choose the second, ${8 \choose 4}$ for the third, and ${4 \choose 4}=1$ for the fourth. Since the committees are indistinguishable, we need to divide the product by $4!$. Thus the $16$ people can be grouped in
 \[\frac{1}{4!}{16 \choose 4}{12 \choose 4}{8 \choose 4}{4 \choose 4}=\frac{16!}{(4!)^5}\]ways.

 In each committee, there are $4 \cdot 3=12$ ways to choose the chairperson and secretary, so $12^4$ ways for all $4$ committees. Therefore, there are
 \[\frac{16!}{(4!)^5}12^4\]total possibilities.

 Since $16!$ contains $6$ factors of $3$, $(4!)^5$ contains $5$, and $12^4$ contains $4$, $r=6-5+4=\boxed{\textbf{(A) }5}$.

 



\textbf{Solution 2 (Multinomial Coefficients)}

There are $\binom{16}{4,4,4,4}$ ways to choose the 4 committees. You have to divide by another 4! since the order of the committees does not matter. Furthermore, in each committee, you can have $4 \cdot 3$ ways to choose chairperson and secretary. Hence a total of $\lfloor{\frac{16}{3}\rfloor}+\lfloor{\frac{16}{9}\rfloor}+4-5=5.$




\textbf{Solution 3}

There will be $16$ ways to pick the chairperson of the first committee, then $15$ to pick the secretary, and lastly ${14 \choose 2}$ ways to pick the other two members of the first committee. Similarly, we can complete the rest of the terms as follows:
 \[\frac{(16)(15){14 \choose 2}(12)(11){10 \choose 2}(8)(7){6\choose 2}(4)(3){2\choose 2}}{4!}\]We notice the numerator has at most $3^6$, and the denominator has just $3$. Thus, the value of $r$ in question is $\boxed{\textbf{(A)}\ 5}$.

 



\textbf{Solution 4}

We can convert the question into arranging people in a line of length $16$.

 First, arrange all $16$ people in a line. There are $16!$ possible ways to do this. Then, divide the line into $4$ groups of $4$, which represent the committees. Since the committees are indistinguishable, we divide by the number of ways to permute the $4$ groups, $4!$.  

 Within each group of $4$, the first two positions in the line are assigned as the chairperson and the secretary, while the remaining two are ordinary members. Since the roles of chairperson and secretary are distinct, but the order of the remaining two members does not matter, we divide by $2$ for each group. Thus, we divide by $2^4$ to account for all $4$ committees.  

 The total number of ways to assign people into committees with roles is given by:
 \[\frac{16!}{4! \cdot 2^4}.\]  (This is equivalent to Solution 1's $\frac{16!}{(4!)^5} \cdot 12^4$.)

 To find $r$, we use Legendre's Formula to count the total power of $3$ in $16!$:  
 \[\lfloor \frac{16}{3} \rfloor + \lfloor \frac{16}{9} \rfloor = 5 + 1 = 6.\]  Each $4!$ contributes one factor of $3$, and since we divide by $4!$ a total of $1$ times, we subtract $1$ from $6$:  
 \[6 - 1 = 5.\]  Thus, the value of $r$ is $\boxed{5}$.~ryk

 


	\item (\textit{AMC 10 10B 2024}) The Fibonacci numbers are defined by $F_1 = 1, F_2 = 1,$ and $F_n = F_{n-1} + F_{n-2}$ for $n \geq 3.$ What is 
 \[{\frac{F_2}{F_1}} + {\frac{F_4}{F_2}} + {\frac{F_6}{F_3}} + ... + {\frac{F_{20}}{F_{10}}}?\]
$\textbf{(A) } 318 \qquad\textbf{(B) } 319 \qquad\textbf{(C) } 320 \qquad\textbf{(D) } 321 \qquad\textbf{(E) } 322$

 \textbf{Answer: }B



\textbf{Solution 1 (Brute Force)}

The first $20$ terms are $F_n = 1, 1, 2, 3, 5, 8, 13, 21, 34, 55, 89, 144, 233, 377, 610, 987, 1597, 2584, 4181, 6765$

 So the answer is $1 +  3 + 4 + 7 + 11 + 18 + 29 + 47 + 76 + 123 =  \boxed{(B) 319}$. 

 
            \begin{itemize}
        \item Do not do this unless you have no other option or no time to find a smarter way
            \end{itemize}
            


	\item (\textit{AMC 10 10B 2024}) Let

 \[P(m)=\frac{m}{2}+\frac{m^2}{4}+\frac{m^4}{8}+\frac{m^8}{8}\]
How many of the values $P(2022)$, $P(2023)$, $P(2024)$, and $P(2025)$ are integers?


 $\textbf{(A) } 0 \qquad\textbf{(B) } 1 \qquad\textbf{(C) } 2 \qquad\textbf{(D) } 3 \qquad\textbf{(E) } 4$

 \textbf{Answer: }E



\textbf{Solution 1 (The simplest way)}

First, we know that $P(2022)$ and $P(2024)$ must be integers since they are both divisible by $2$. 

 Then let’s consider the remaining two numbers. Since they are not divisible by $2$, the result of the first term must be a certain number $+\frac{1}{2}$. 

 The divisibility rule for $4$ states that the last two digits of the number must be divisible by $4$, and we know that the last two digits of ${2023}^2$ are the last two digits of ${23}^2$, which is $29$. And $29$ is $1$ greater than a multiple of $4$. Therefore, the result of the second term must be $\frac{1}{4}$. Similarly, the remaining two terms must each be $\frac{1}{8}$.

 Their sum is $\frac{1}{2} + \frac{1}{4} + \frac{1}{8} + \frac{1}{8} = 1$, so $P(2023)$ and $P(2025)$ are also integers.

 Therefore, the answer is $\boxed{\textbf{(E) }4}$.

 



\textbf{Solution 2 (Specific)}

Take everything modulo 8 and re-write the entire fraction with denominator 8. This means that we're going to transform the fraction as follows : 
 \[P(m)=\frac{m}{2}+\frac{m^2}{4}+\frac{m^4}{8}+\frac{m^8}{8}\]becomes 
 \[P(m)=\frac{4m + 2m^2 + m^4 + m^8}{8}\]And in order for $P(m)$ to be an integer, it's important to note that $4m + 2m^2 + m^4 + m^8$ must be congruent to 0 modulo 8. Moreover, we know that $2022 \equiv 6 \pmod 8, 2023 \equiv 7 \pmod 8, 2024 \equiv 0 \pmod 8, 2025 \equiv 1 \pmod 8$. We can verify it by taking everything modulo 8 : 

 If $m = 2022$, then $4(6) + 2(4) + 0 + 0 = 24 + 8 = 32 \equiv 0 \pmod 8$ -> TRUEIf $m = 2023$, then $4(7) + 2(1) + 1 + 1 = 28 + 2 + 1 + 1 = 32 \equiv 0 \pmod 8$ -> TRUEIf $m = 2024$, then it is obvious that the entire expression is divisible by 8. Therefore, it is true. If $m = 2025$, then $2025 \equiv 1 \pmod 8$. Therefore, $4(1) + 2(1) + 1 + 1 = 8 \equiv 0 \pmod 8$ -> TRUETherefore, there are $\boxed{\textbf{(E) }4}$ possible values. 

 Addendum for certain China test papers : Note that $2026 \equiv 2 \pmod 8$. Therefore, taking everything modulo 8, whilst still maintaining the original expression, gives $4(2) + 2(4) + 0 + 0 = 16 \equiv 0 \pmod 8$. This is true. 

 



\textbf{Solution 3 (Factoring)}

We can rewrite the expression as 
 \[\frac{1}{8}\cdot m \cdot (m^7+m^3+2m+4).\] If $m$ is even, then $m$ gives a factor of $2$ and $m^7+m^3+2m+4$ will give a factor of $4,$ so the result will be an integer. If $m$ is odd, then notice that for any $m$ we have $m^2\equiv 1 \pmod{8}.$ Then $m^7+m^3+2m+4\equiv m+m+2m+4 \equiv 4(m+1) \equiv 0 \pmod{8}.$ So any integer $m$ will result in an integer, meaning the answer is $\boxed{\textbf{(E) }4}$.

 -nevergonnagiveup

 



\textbf{Solution 4}

Taking a closer look at the terms, we notice that each term builds off of the previous one. There is $\frac{m}{2}$, $\frac{m^{2}}{4}$, which is equal to $\left(\frac{m}{2}\right)^{2}$, $\frac{m^{4}}{8}$, which is equal to $2\left(\frac{m^{2}}{4}\right)^{2}$, and $\frac{m^{8}}{8}$, which is equal to $8\left(\frac{m^{4}}{8}\right)^{2}$. Ultimately, this means that there are only two cases that we need to check for: the case in which $m$ is even and the case in which $m$ is odd.

 If $m$ is even, $\frac{m}{2}$ will be an integer, which means the rest of the terms will be an integer. This means in the problem, $P(2022)$ and $P(2024)$ will yield an integer result. (For certain Chinese versions, this also proves that $P(2026)$ is an integer.)

 If $m$ is odd, $\frac{m}{2}$ will result in $\frac{1}{2}$. Building off of each term will give you $\frac{1}{4}$, $\frac{1}{8}$, and $\frac{1}{8}$, and summing those up will grant $1$, an integer. This means in the problem, $P(2023)$ and $P(2025)$ will also result in integer answers.

 In general, all integer $m$ will make $P(m)$ give an integer answer, but for this question, this will get us to $\boxed{\textbf{(E) }4}$ integer values ($\boxed{\textbf{(E) }5}$ for certain Chinese versions of the test).

 



\textbf{Solution 5 (sol 2 but fast)}

Obviously $P(2022)$ and $P(2024)$ both work. $2023 \equiv -1 \pmod 8$, and $2025 \equiv 1 \pmod 8$. Restating the problem, we have the equation equivalent to $\frac{4m+2m^2 + m^4 + m^8}{8}$. Plugging in $m=-1$, we have $-4+2+1+1 \equiv 0 \pmod 8$, so $P(2023)$ works. Plugging in $m=1$, we have $4+2+1+1 \equiv 0 \pmod 8$ as well, so $P(2025)$ also works. Thus, the answer is $\boxed{E}$.

 



\textbf{Solution 6}

Note that $2022 \equiv  0\pmod{2}$, $2\pmod{4}$, and $-2\pmod{8}$. 

 Thus, when $m=2022$, $P(m) = \frac{2022}{2} + \frac{2022^2}{4} + \frac{2022^4}{8} + \frac{2022^8}{8} \equiv \frac{0}{2} + \frac{2^2}{4} + \frac{(-2)^4}{8} + \frac{(-2)^8}{8} \equiv \frac{0}{8} + \frac{8}{8} + \frac{16}{8} + \frac{256}{8} \equiv \frac{280}{8} \equiv 0\pmod{8}$ meaning $P(2022)|8$ and thus is an integer. 

 Repeating this process for each of $m=2023, 2024,$ and $2025$ shows that they're all integers, giving a final answer of $\boxed{\textbf{(E) 4}}$.

 \textbf{Remark} On certain versions of the AMC in China, the problem was restated as follows: 

 Let
 \[P(m)=\frac{m}{2}+\frac{m^2}{4}+\frac{m^4}{8}+\frac{m^8}{8}\]How many of the values $P(2022)$, $P(2023)$, $P(2024)$, $P(2025),$ and $P(2026)$ are integers?$\textbf{(A) } 1 \qquad\textbf{(B) } 2 \qquad\textbf{(C) } 3 \qquad\textbf{(D) } 4 \qquad\textbf{(E) } 5$

 By identical reasoning, each term of $P$ is an integer, since $2026$ is even. 

 Therefore, the answer is $\boxed{\textbf{(E) }5}$.

 \textbf{Not So Great Guessing Method} Note that this is Question $24$ of the $2024$ AMC $10$B. It follows that the answer should be something to do with 2024. Option $\boxed{\textbf{(E) }4}$ is the closest one with the last digit of $4$. Therefore, if you have no way to solve this and have no clue why you decided to take the AMC, use this method.(Like me!)

 


	\item (\textit{AMC 10 10B 2024}) Each of $27$ bricks (right rectangular prisms) has dimensions $a \times b \times c$, where $a$, $b$, and $c$ are pairwise relatively prime positive integers. These bricks are arranged to form a $3 \times 3 \times 3$ block, as shown on the left below. A $28$th brick with the same dimensions is introduced, and these bricks are reconfigured into a $2 \times 2 \times 7$ block, shown on the right. The new block is $1$ unit taller, $1$ unit wider, and $1$ unit deeper than the old one. What is $a + b + c$?


 
                    \begin{figure}[H]
                        \centering
                        \includegraphics[width=0.5\textwidth]{Problem/AMC_10_10B_2024_25_Problem_Figure_1.png}
                    \end{figure}
                    
$\textbf{(A) } 88 \qquad\textbf{(B) } 89 \qquad\textbf{(C) } 90 \qquad\textbf{(D) } 91 \qquad\textbf{(E) } 92$

 \textbf{Answer: }E



\textbf{Solution 1 (Less than 60 seconds)}

The $3$x$3$x$3$ block has side lengths of $3a, 3b, 3c$. The $2$x$2$x$7$ block has side lengths of $2b, 2c, 7a$.

 We can create the following system of equations, knowing that the new block has $1$ unit taller, deeper, and wider than the original:
 \[3a+1 = 2b\]
 \[3b+1=2c\]
 \[3c+1=7a\]

 Adding all the equations together, we get $b+c+3 = 4a$. Adding $a-3$ to both sides, we get $a+b+c = 5a-3$. The question states that $a,b,c$ are all relatively prime positive integers. Therefore, our answer must be congruent to $2 \pmod{5}$. The only answer choice satisfying this is $\boxed{(E)\ 92}$.~lprado

 



\textbf{Solution 2 (Solution 1 but longer)}

We will define the equations the same as solution 1.
 \[3a+1 = 2b\]
 \[3b+1=2c\]
 \[3c+1=7a\]Solve equation 2 for c and substitute that value in for equation 3, giving us
 \[3a+1 = 2b\]
 \[\frac{3b+1}{2}=c\]
 \[\frac{3 \cdot (3b+1)}{2}+1=7a\]

 Multiply 14 to the first equation and rearrange to get
 \[42a = 28b-14\]and multiply the third by 2 and rearrange to get
 \[27b+15 = 42a\]Solve for b to get $b = 29$, substitute into equation 1 from the original to get $a = 19$, and lastly, substitute a into original equation 2 to get $c = 44$.Thus, $a+b+c = 19+29+44 = \boxed{E(92)}$.~Failure.net

 



\textbf{Solution 3}

We will define the equations the same as solution 1 and 2.
 \[3a+1 = 2b\]
 \[3b+1=2c\]
 \[3c+1=7a\]Multiply 1st equation by 1.5 to get
 \[4.5a+1.5=3b\]Add 1 to get
 \[4.5a+2.5=3b+1\]Substitute to get
 \[4.5a+2.5=2c\]Multiply this one by 1.5 as well, getting
 \[6.75a+3.75=3c\]Adding 1 again to get
 \[6.75a+4.75=3c+1\]Replace 
 \[3c+1\] with 
 \[7a\] getting 
 \[6.75a+4.75=7a\]Solve for 
 \[a\] getting 
 \[0.25a=4.75\] 
 \[a=19\]

 Solve for b: 
 \[3(19)+1=2b\] 
 \[57+1=2b\] 
 \[b=29\]

 Solve for c: 
 \[3(29)+1=2c\] 
 \[87+1=2c\] 
 \[c=44\]Thus, $a+b+c = 19+29+44 = \boxed{E(92)}$




\textbf{Solution 4 (No Fractions)}

As stated in solutions 1, 2, 3, the equations are as follows:(1) 
 \[3a+1=2b\](2) 
 \[3b+1=2c\](3) 
 \[3c+1=7a\]Multiply equation (2) by 2 to get equation (4):(4) 
 \[6b+2=4c\]and substitute (1) into (4):(5) 
 \[9a+5=4c\](3) 
 \[7a-1=3c\]Multiplying (5) by 3 and (3) by 4 gives:
 \[3(9a+5)=4(7a-1)\]from which we get $a=19$. Substituting $a$ into (1) and (3) yields $b=29$ and $c=44$. Therefore, $a+b+c = \boxed{E(92)}$



\end{enumerate}